% ===================================================================
% Paper M1 — Four axioms and their minimal geometry
% Base paper of the Master Equation Framework series.  Draft v1, 7 Nov 2026.
% Compile: xelatex (three passes).  Certificates: M1_lean/*.lean, Lean 4.15.0 core.
% ===================================================================
\documentclass[11pt,a4paper]{article}
\usepackage{amsmath,amssymb,amsthm}
\usepackage[margin=25mm]{geometry}
\usepackage{enumitem}
\usepackage{booktabs}
\usepackage{fontspec}
\setmonofont{DejaVu Sans Mono}[Scale=0.78]
\usepackage{listings}
\lstset{basicstyle=\ttfamily\footnotesize,breaklines=true,columns=fullflexible,keepspaces=true,frame=single,framesep=4pt,xleftmargin=2pt,extendedchars=true}
\usepackage[hidelinks]{hyperref}
%%=====================================================================
%% mef_macros.tex
%% Canonical macro set for the Master Equation Framework (MEF) series.
%% Version: 1.0 (April 2026)
%%
%% Single source of truth for symbol notation across MEF papers.
%% Derived from main.tex (thesis preamble) plus selected additions
%% from Paper XVI (number sets and rigour-classification marks).
%%
%% Usage: place %%=====================================================================
%% mef_macros.tex
%% Canonical macro set for the Master Equation Framework (MEF) series.
%% Version: 1.0 (April 2026)
%%
%% Single source of truth for symbol notation across MEF papers.
%% Derived from main.tex (thesis preamble) plus selected additions
%% from Paper XVI (number sets and rigour-classification marks).
%%
%% Usage: place %%=====================================================================
%% mef_macros.tex
%% Canonical macro set for the Master Equation Framework (MEF) series.
%% Version: 1.0 (April 2026)
%%
%% Single source of truth for symbol notation across MEF papers.
%% Derived from main.tex (thesis preamble) plus selected additions
%% from Paper XVI (number sets and rigour-classification marks).
%%
%% Usage: place \input{mef_macros.tex} in the preamble of every MEF
%% paper after the package loads but before \begin{document}.  Drift
%% between papers is prevented by editing only this file.
%%
%% Companion file: mef_field_content_preamble.tex (Section 2 reader-
%% orientation block).
%%=====================================================================

%% ---- Quantum Coherence Dilaton and its constituents ----
\newcommand{\QCD}{\Phi_{\mathrm{QC}}}           % Quantum Coherence Dilaton (canonical)
\newcommand{\consciousness}{\omega}             % consciousness coherence field
\newcommand{\zetaval}{12\pi}                    % quantum coherence coupling = chi(K_8) * pi

%% ---- Manifolds ----
\newcommand{\Mtwelve}{\mathcal{M}_{12}}         % 12D parent manifold
\newcommand{\Keight}{K_8}                       % 8D internal manifold
\newcommand{\CPtwo}{\mathbb{CP}^2}              % complex projective plane
\newcommand{\Stwo}{S^2}                         % 2-sphere
\newcommand{\Ttwo}{T^2}                         % 2-torus

%% ---- Discrete groups ----
\newcommand{\Ztwo}{\mathbb{Z}_2}                % pillowcase orbifold group
\newcommand{\Zfour}{\mathbb{Z}_4}

%% ---- Spin^c and bundles ----
\newcommand{\Spinc}{\text{Spin}^c}              % Spin^c structure
\newcommand{\Oline}[1]{\mathcal{O}(#1)}         % line bundle of degree #1

%% ---- Number fields (NB: \Q = mathbb Q is RATIONALS, not the dilaton; \QCD = \mathcal{Q} is the dilaton) ----
\newcommand{\C}{\mathbb{C}}
\newcommand{\R}{\mathbb{R}}
\newcommand{\Z}{\mathbb{Z}}
\newcommand{\Q}{\mathbb{Q}}

%% ---- Effective couplings, scales, moduli ----
\newcommand{\Veff}{V_{\text{eff}}}
\newcommand{\VCKM}{V_{\text{CKM}}}
\newcommand{\lambdaC}{\lambda_C}
\newcommand{\tautwo}{\tau_2}                    % stabilised aspect modulus
\newcommand{\thetaW}{\theta_W}                  % Weinberg angle
\newcommand{\Rw}{R_w}                           % S^2 radius (warped fibre)
\newcommand{\Rc}{R_c}                           % CP^2 radius (Calabi-like)
\newcommand{\Rfibre}{\mathcal{R}_{\text{fibre}}}
\newcommand{\Cfibre}{C_{\mathrm{fibre}}}
\newcommand{\Rsix}{\mathcal{R}_6}
\newcommand{\Fsq}{\mathcal{F}^2}
\newcommand{\centcorr}{\mathcal{C}_{32}}

%% ---- Spin^c holonomy phases ----
\newcommand{\Phiomega}{\phi_\consciousness}     % omega-axis holonomy
\newcommand{\Phipsi}{\phi_\psi}                 % psi-axis holonomy

%% ---- Arithmetic ----
\newcommand{\dodd}{d_{\mathrm{odd}}}            % odd divisor function

%% ---- Operators ----
\newcommand{\Tr}{\mathrm{Tr}}
\newcommand{\Vol}{\mathrm{Vol}}
\newcommand{\vol}{\mathrm{vol}}

%% ---- Rigour-classification marks ----
\newcommand{\rigmark}{\hfill\textbf{[Rigorous]}}
\newcommand{\dermark}{\hfill\textbf{[Derived]}}
\newcommand{\motmark}{\hfill\textbf{[Motivated]}}
\newcommand{\gapmark}{\hfill\textbf{[Gap]}}
\newcommand{\rigder}{\hfill\textbf{[Rigorous / Derived]}}
\newcommand{\dergap}{\hfill\textbf{[Derived / Gap]}}

%% ---- End of mef_macros.tex ----
 in the preamble of every MEF
%% paper after the package loads but before \begin{document}.  Drift
%% between papers is prevented by editing only this file.
%%
%% Companion file: mef_field_content_preamble.tex (Section 2 reader-
%% orientation block).
%%=====================================================================

%% ---- Quantum Coherence Dilaton and its constituents ----
\newcommand{\QCD}{\Phi_{\mathrm{QC}}}           % Quantum Coherence Dilaton (canonical)
\newcommand{\consciousness}{\omega}             % consciousness coherence field
\newcommand{\zetaval}{12\pi}                    % quantum coherence coupling = chi(K_8) * pi

%% ---- Manifolds ----
\newcommand{\Mtwelve}{\mathcal{M}_{12}}         % 12D parent manifold
\newcommand{\Keight}{K_8}                       % 8D internal manifold
\newcommand{\CPtwo}{\mathbb{CP}^2}              % complex projective plane
\newcommand{\Stwo}{S^2}                         % 2-sphere
\newcommand{\Ttwo}{T^2}                         % 2-torus

%% ---- Discrete groups ----
\newcommand{\Ztwo}{\mathbb{Z}_2}                % pillowcase orbifold group
\newcommand{\Zfour}{\mathbb{Z}_4}

%% ---- Spin^c and bundles ----
\newcommand{\Spinc}{\text{Spin}^c}              % Spin^c structure
\newcommand{\Oline}[1]{\mathcal{O}(#1)}         % line bundle of degree #1

%% ---- Number fields (NB: \Q = mathbb Q is RATIONALS, not the dilaton; \QCD = \mathcal{Q} is the dilaton) ----
\newcommand{\C}{\mathbb{C}}
\newcommand{\R}{\mathbb{R}}
\newcommand{\Z}{\mathbb{Z}}
\newcommand{\Q}{\mathbb{Q}}

%% ---- Effective couplings, scales, moduli ----
\newcommand{\Veff}{V_{\text{eff}}}
\newcommand{\VCKM}{V_{\text{CKM}}}
\newcommand{\lambdaC}{\lambda_C}
\newcommand{\tautwo}{\tau_2}                    % stabilised aspect modulus
\newcommand{\thetaW}{\theta_W}                  % Weinberg angle
\newcommand{\Rw}{R_w}                           % S^2 radius (warped fibre)
\newcommand{\Rc}{R_c}                           % CP^2 radius (Calabi-like)
\newcommand{\Rfibre}{\mathcal{R}_{\text{fibre}}}
\newcommand{\Cfibre}{C_{\mathrm{fibre}}}
\newcommand{\Rsix}{\mathcal{R}_6}
\newcommand{\Fsq}{\mathcal{F}^2}
\newcommand{\centcorr}{\mathcal{C}_{32}}

%% ---- Spin^c holonomy phases ----
\newcommand{\Phiomega}{\phi_\consciousness}     % omega-axis holonomy
\newcommand{\Phipsi}{\phi_\psi}                 % psi-axis holonomy

%% ---- Arithmetic ----
\newcommand{\dodd}{d_{\mathrm{odd}}}            % odd divisor function

%% ---- Operators ----
\newcommand{\Tr}{\mathrm{Tr}}
\newcommand{\Vol}{\mathrm{Vol}}
\newcommand{\vol}{\mathrm{vol}}

%% ---- Rigour-classification marks ----
\newcommand{\rigmark}{\hfill\textbf{[Rigorous]}}
\newcommand{\dermark}{\hfill\textbf{[Derived]}}
\newcommand{\motmark}{\hfill\textbf{[Motivated]}}
\newcommand{\gapmark}{\hfill\textbf{[Gap]}}
\newcommand{\rigder}{\hfill\textbf{[Rigorous / Derived]}}
\newcommand{\dergap}{\hfill\textbf{[Derived / Gap]}}

%% ---- End of mef_macros.tex ----
 in the preamble of every MEF
%% paper after the package loads but before \begin{document}.  Drift
%% between papers is prevented by editing only this file.
%%
%% Companion file: mef_field_content_preamble.tex (Section 2 reader-
%% orientation block).
%%=====================================================================

%% ---- Quantum Coherence Dilaton and its constituents ----
\newcommand{\QCD}{\Phi_{\mathrm{QC}}}           % Quantum Coherence Dilaton (canonical)
\newcommand{\consciousness}{\omega}             % consciousness coherence field
\newcommand{\zetaval}{12\pi}                    % quantum coherence coupling = chi(K_8) * pi

%% ---- Manifolds ----
\newcommand{\Mtwelve}{\mathcal{M}_{12}}         % 12D parent manifold
\newcommand{\Keight}{K_8}                       % 8D internal manifold
\newcommand{\CPtwo}{\mathbb{CP}^2}              % complex projective plane
\newcommand{\Stwo}{S^2}                         % 2-sphere
\newcommand{\Ttwo}{T^2}                         % 2-torus

%% ---- Discrete groups ----
\newcommand{\Ztwo}{\mathbb{Z}_2}                % pillowcase orbifold group
\newcommand{\Zfour}{\mathbb{Z}_4}

%% ---- Spin^c and bundles ----
\newcommand{\Spinc}{\text{Spin}^c}              % Spin^c structure
\newcommand{\Oline}[1]{\mathcal{O}(#1)}         % line bundle of degree #1

%% ---- Number fields (NB: \Q = mathbb Q is RATIONALS, not the dilaton; \QCD = \mathcal{Q} is the dilaton) ----
\newcommand{\C}{\mathbb{C}}
\newcommand{\R}{\mathbb{R}}
\newcommand{\Z}{\mathbb{Z}}
\newcommand{\Q}{\mathbb{Q}}

%% ---- Effective couplings, scales, moduli ----
\newcommand{\Veff}{V_{\text{eff}}}
\newcommand{\VCKM}{V_{\text{CKM}}}
\newcommand{\lambdaC}{\lambda_C}
\newcommand{\tautwo}{\tau_2}                    % stabilised aspect modulus
\newcommand{\thetaW}{\theta_W}                  % Weinberg angle
\newcommand{\Rw}{R_w}                           % S^2 radius (warped fibre)
\newcommand{\Rc}{R_c}                           % CP^2 radius (Calabi-like)
\newcommand{\Rfibre}{\mathcal{R}_{\text{fibre}}}
\newcommand{\Cfibre}{C_{\mathrm{fibre}}}
\newcommand{\Rsix}{\mathcal{R}_6}
\newcommand{\Fsq}{\mathcal{F}^2}
\newcommand{\centcorr}{\mathcal{C}_{32}}

%% ---- Spin^c holonomy phases ----
\newcommand{\Phiomega}{\phi_\consciousness}     % omega-axis holonomy
\newcommand{\Phipsi}{\phi_\psi}                 % psi-axis holonomy

%% ---- Arithmetic ----
\newcommand{\dodd}{d_{\mathrm{odd}}}            % odd divisor function

%% ---- Operators ----
\newcommand{\Tr}{\mathrm{Tr}}
\newcommand{\Vol}{\mathrm{Vol}}
\newcommand{\vol}{\mathrm{vol}}

%% ---- Rigour-classification marks ----
\newcommand{\rigmark}{\hfill\textbf{[Rigorous]}}
\newcommand{\dermark}{\hfill\textbf{[Derived]}}
\newcommand{\motmark}{\hfill\textbf{[Motivated]}}
\newcommand{\gapmark}{\hfill\textbf{[Gap]}}
\newcommand{\rigder}{\hfill\textbf{[Rigorous / Derived]}}
\newcommand{\dergap}{\hfill\textbf{[Derived / Gap]}}

%% ---- End of mef_macros.tex ----


\theoremstyle{plain}
\newtheorem{theorem}{Theorem}[section]
\newtheorem{proposition}[theorem]{Proposition}
\newtheorem{lemma}[theorem]{Lemma}
\newtheorem{corollary}[theorem]{Corollary}
\theoremstyle{definition}
\newtheorem{definition}[theorem]{Definition}
\newtheorem{axiom}{Axiom}
\theoremstyle{remark}
\newtheorem{remark}[theorem]{Remark}

\newcommand{\Hquat}{\mathbb{H}}
\newcommand{\ImH}{\operatorname{Im}\Hquat}
\newcommand{\Ad}{\operatorname{Ad}}
\newcommand{\Fix}{\operatorname{Fix}}
\newcommand{\Sol}{\mathrm{Sol}}
\newcommand{\rigR}{[\textup{R}]}
\newcommand{\rigD}{[\textup{D}]}
\newcommand{\rigM}{[\textup{M}]}
\newcommand{\rigG}{[\textup{G}]}

\title{The Master Equation from four axioms: the field equation of the minimal twelve-dimensional geometry}
\author{Dhiren J. Master}
\date{9 November 2026 (v1.0, canon)}

\begin{document}
\maketitle

\begin{abstract}
Einstein's field equations relate the curvature of a space to the matter in it and to the energy of empty space, and in a century of use they have described gravity and left the other forces, the existence and structure of matter, and the direction of time to be supplied from outside. The first paper of this series derived, from four sentences a child can verify, a unique twelve-dimensional geometry. We show here that the same four sentences, with the same principle of minimality, derive a field equation on that geometry, and that the equation is Einstein's with one additional term. Two of the axioms, closure and order, give Einstein's equation with its cosmological and matter terms. The third, that the universe grows and decays as a function of itself, gives the additional term: the exponential of growth and decay has two sheets, the sheets are one field $\psi$, a section of the Clifford module of the geometry, with shadow $\bar\psi$, their even combination $\omega=\psi+\bar\psi$ and the unique lowest-order even scalar vanishing both where the sheets coincide and where growth opposes decay, $\QCD=\omega\lvert\psi-\bar\psi\rvert^{2}$, enter the action through one coupling $\zeta$, and the variation gives $G_{MN}=8\pi G_{\mathrm{eff}}T_{MN}+\zeta T^{\QCD}_{MN}+\Lambda_{\mathrm{eff}}g_{MN}$. The fourth axiom, the drive, makes the two sheets mirror images relative to the flow of time rather than one sheet written backwards, so that the term is physical and the arrow of time and the distinctness of matter from its shadow are the same fact; no Lorentzian structure is assumed. The coupling is derived from the restriction of the equation to the base of the geometry as the exponent by which the warp falls across a half-period, $\zeta=\chi(\Keight)\pi=12\pi$. The derivation is minimal rather than unique: at four named points the smallest option the axioms admit is taken, and each is tagged. The field $\omega$ vanishes at the locked phase by the second axiom and at the uncorrelated phase by the fourth, and the coupling is largest at the sheet angle $\arctan\sqrt2$.
\end{abstract}

\section{What is observed, and what this paper derives}\label{sec:intro}

The purpose of this paper is to derive the field equation of the Master Equation programme from the four axioms of M1, on the geometry M1 showed those axioms select, and to show that the equation is Einstein's with one additional term, produced by the one axiom Einstein's theory does not contain.

\subsection{The result in plain words}\label{sec:plain}

M1 began from four observations, taken as axioms in the simplest words that carry them, and we restate them as M1 did.

\begin{axiom}[closure]\label{ax:A1}
The universe closes on itself, so that it can be aware of itself.
\end{axiom}

\begin{axiom}[growth and decay]\label{ax:A2}
The universe grows and decays as a function of itself.
\end{axiom}

\begin{axiom}[order]\label{ax:A3}
The order of events matters, so that each instance is unique and cannot be foretold.
\end{axiom}

\begin{axiom}[drive]\label{ax:A4}
The universe is in constant change; it never stands still.
\end{axiom}

Read as conditions on a shape, with every point at which a word acquires a mathematical meaning isolated as a numbered reading, and with the rule that no unforced structure is added, these four picked out one shape, twelve-dimensional, described in M1. This paper asks what rule of motion the same four sentences impose on that shape.

The answer can be stated before any mathematics. Einstein's equation says that the curvature of space is produced by two things: the matter in it, and the energy that empty space has of its own. Our equation says the same, with one more source of curvature added: the difference between growth and decay, weighted by how coherently the two are held together. Written as a formula it is Einstein's equation with one extra term on the right-hand side, and the extra term has one number in front of it, which we compute to be $12\pi$.

Which observation produces which part of the equation is the content of the paper, and it is worth saying at once. That the universe closes on itself is what makes an equation of motion exist at all, and it supplies the energy of empty space; that the order of events matters is what supplies matter, in the form of the spinor fields that carry non-commutativity, together with the forces that act between them. These two observations alone give Einstein's equation, with its matter and its cosmological constant. That the universe grows and decays is the observation Einstein's theory does not contain, and it is the one that produces the additional term: the growth law has two faces, growth and decay, and the field written on those two faces is a new source of curvature. That the universe never stands still is what makes the two faces genuinely two rather than one face read in a mirror; it is the source of the arrow of time, and it is why matter is distinct from its shadow. In the programme that follows this paper, the second observation is also where the masses, the generations and the mixings of matter come from, so that what was missing from a theory of four forces, an account of why there is matter and why it does not run backwards, is supplied by the two observations the earlier theory left out. We do not derive those consequences here; we derive the equation and its coupling.

\subsection{What the mathematics is doing}\label{sec:what-maths}

For a reader meeting this for the first time, the argument has three moves, and each can be said in ordinary words.

The first move is to write down a single number, the \emph{action}, for
every way the shape and the fields on it could possibly be arranged, and
then to pick out the arrangement for which that number is \emph{stationary}.
Stationary means this: if you alter the arrangement by a tiny amount in any
direction --- bend the shape slightly here, raise the field slightly there
--- the number does not change to first order. It is the same idea as finding
the bottom of a valley by looking for the place where the ground is level:
at the lowest point, a small step in any direction neither raises nor
lowers you. For a function of one variable this is the condition that its
derivative be zero; the action is a function of infinitely many variables
(the value of every field at every point), and the stationary condition is
that its derivative with respect to each of them be zero at once. Working
out what those conditions say is a mechanical calculation, and the result
is the field equation. This is how every field equation in physics is
obtained, Einstein's included: the equation is not guessed, it is the
condition that the action be stationary. The first axiom is what makes the
move possible, because a shape that closes on itself has no edge over which
the small alterations could leak away and spoil the count.

The content of the paper is then entirely in \emph{which terms the action
is allowed to contain}. We allow nothing the axioms do not force, and where
the axioms leave a choice we take the smallest option, as M1 did; the paper
names each such choice. There are three kinds of term: the curvature of the
shape itself, the matter fields and the force field that the third axiom
supplies, and the field written on the two faces of the growth law that the
second axiom supplies.

The second move is to identify that last field. The growth law $w=w_{0}e^{\kappa s}$ has two solutions, growth and decay, exchanged by an operation the second axiom asks for. We show that this operation is the reversal of the drive of the fourth axiom together with a quarter-turn of the internal fibre, that it reverses the orientation of the shape, and hence that the two solutions, written as one field $\psi$ and its shadow $\bar\psi$, have opposite handedness. The even combination $\omega=\psi+\bar\psi$ measures how far the two faces agree; the odd combination $\psi-\bar\psi$ measures how far they are held apart. The one scalar built from them that respects the exchange and vanishes both when the faces coincide and when growth exactly opposes decay is $\QCD=\omega\lvert\psi-\bar\psi\rvert^{2}$, and the shape couples to the two faces through it and through nothing else.

The third move is the computation. Varying the action gives Einstein's equation with the one extra term; the extra term is the stress tensor of a scalar field, a standard object, multiplied by the coupling $\zeta$. To find $\zeta$ we restrict the equation to the two-dimensional base of the internal shape, where the growth law appears as a warp that falls off across the base, and we show that across one half of the base the warp falls by exactly $N\pi$, where $N$ is the number of units of flux the shape carries and $\pi$ is the angle missing at each of the four corners of the pillowcase. The number of flux units is the Euler characteristic of the internal shape, twelve, so $\zeta=12\pi$.

\subsection{What this paper proves, and how it is organised}\label{sec:results}

The results, in the order they are proved, are these. The sheets of A2 are one field $\psi$, a section of the Clifford module $\Hquat^{32}$ that A3 forces, with shadow $\bar\psi$ obtained by the flow-reversing involution; $\omega=\psi+\bar\psi$ and $\psi-\bar\psi$ are the Coherent Field and the sheet variable, and $\QCD=\omega\lvert\psi-\bar\psi\rvert^{2}$ is the unique lowest-order even scalar vanishing both where the sheets coincide and where growth opposes decay (Section~\ref{sec:variables}). The involution is the quarter-turn flow reversal $\tau$ with $\tau^{2}$ the fold, it reverses orientation, and $\bar\psi$ has the opposite handedness to $\psi$ (Section~\ref{sec:handedness}). The action on $M_{12}$ that the axioms and the standing minimality principle admit has one gravitational term, the induced connection of the $\Spinc$ line, and the dilaton coupled through $\QCD$ alone; its variation is the Master Equation, Einstein's equation with the one additional term, and the dilaton's equation is $\zeta\Box\QCD=V'$ (Section~\ref{sec:equation}). The restriction of the equation to the base is the two-dimensional system in which the warp falls by $N\pi$ across a half-period, $N=\chi(\Keight)=12$, so $\zeta=12\pi$ (Section~\ref{sec:warp}). The Coherent Field vanishes at the locked phase by A2 and at the uncorrelated phase by A4, the coupling is largest at $\arctan\sqrt2$, and $\omega$ keeps one sign on the solutions (Section~\ref{sec:omega-shape}).

The paper is organised as follows. Section~\ref{sec:variables} defines the variables, states the fifth reading of the series, that the metric is positive-definite and time is the drive's orbit parameter, and states the standing principle of minimality. Section~\ref{sec:handedness} proves the handedness theorem. Section~\ref{sec:equation} writes the action and derives the equation, saying for each term which axiom produces it. Section~\ref{sec:warp} derives the value of $\zeta$. Section~\ref{sec:omega-shape} derives the shape of $\omega$. Section~\ref{sec:outlook} says what the equation dictates for the physics and what remains open. As in M1, this introduction and the opening of each section are intended to be followed by a first-year student; from Section~\ref{sec:handedness} the register is that of a formal paper, every load-bearing step is proved or reduced to a cited classical result or to a stated reading, the smallest case is worked by hand, and the finite and arithmetic parts are certified in Lean (Appendix~\ref{app:lean}). Rigour is tagged \rigR, \rigD, \rigM, \rigG\ as in M1. The deflationary reading of the paper is stated where it belongs: the equation is derived from the action, and the action is fixed by the axioms together with four minimal choices, each named (Sections~\ref{sec:summary3} and~\ref{sec:summary4}); a reader who rejects minimality as a principle has a family of actions, not an equation. Every symbol of M1 keeps its meaning; the symbols introduced here are tabulated in Appendix~\ref{app:labels}.

%%=====================================================================
%% M2 — WP0 definitions sheet. SURGICAL FRAGMENT (insertion only; no
%% preamble; symbols per mef_macros.tex and the M1 preamble macros
%% \Hquat, \Ad, \rigR, \rigD, \rigM, \rigG). UK English. "We" throughout.
%% Architect draft, 8 Nov 2026 — every choice that is the architect's
%% rather than a PI ruling or an M1 result is marked in a %% ARCH comment.
%% Intended location: M2 Section 2 (after the recall of A1–A4 and M1).
%%=====================================================================

\section{The variables}\label{sec:variables}

The geometry is fixed by M1: $M_{12}=(S^{1}\times\Sigma)\times\Keight$ with
$\Sigma=S^{3}$ and $\Keight=[(\CPtwo\times\Stwo)\times_{w}(\Ttwo/\Ztwo)]^{\Spinc}$,
the fibre $\Ttwo$ carrying the involution of A2 (the fold on the pillowcase; refined in Section~\ref{sec:handedness}) and the
Fibonacci drive $M_{0}$ of A4, with corners $P_{1},\dots,P_{4}$ labelled as in
M1 Appendix~B. We keep every convention of M1 unchanged; in particular the
fibre coordinates are $(\theta_{\psi},\theta_{\omega})$ with half-periods
$(L_{\psi},L_{\omega})$. In this paper $\psi$ and $\omega$ also name
\emph{fields}. The two uses are distinguished by position: a subscript
$\theta_{\psi},\theta_{\omega},L_{\psi},L_{\omega},k_{\omega}$ is a coordinate
label inherited from M1; a bare $\psi$, $\bar\psi$ or $\omega$ is a field
defined in this section. The cumulative table of symbols is Appendix~B.

\subsection{The fifth reading}\label{sec:readings}

M1 numbered four readings by which the four sentences became conditions,
and its abstract already describes the geometry as compact and Riemannian
without making that a numbered reading. This paper makes it one, and
attaches it to every result that depends on it.

\begin{definition}[the metric signs no direction]\label{def:riemannian}
The metric of $M_{12}$ is positive-definite. Time is not a direction the
metric distinguishes by a sign; it is the orbit parameter $s$ of the flow of
A4 on $T\times F$. A derivative ``along time'' is a derivative along that
flow, and every statement below that reads as a statement about time is a
statement along the flow.
\end{definition}

The reading is the geometric form of a sentence as plain as M1's four: the
universe does not label its directions. It is also the reading under which
A1 is unproblematic, since a compact closed manifold with a signed direction
would contain closed time-like curves, and under which the spectra of the
operators of the programme are real and discrete. What singles out a
direction is the drive, and Section~\ref{sec:handedness} shows that the drive
singles it out with an orientation. \rigD
%% PI ruling Q8 (8 Nov 2026). Register: fifth numbered reading of the series.

\paragraph{The standing principle.}
M1 chose among geometries by a minimality principle (M1 Definition~3.2),
and every paper of the series applies the same principle to every datum
the axioms leave open: terms in an action, connections, couplings,
readings. Where a choice exists we take the minimal option and tag the
result \rigD\ with minimality as its hypothesis, on the same footing as
M1's own minimality results; and where the minimal option fails --- a
contradiction, an empty solution set, a conflict with a result already
established --- we say so where it happens rather than enlarge the choice
silently. The present paper applies the principle three times
(Section~\ref{sec:action}), and names each instance.
%% PI, 8 Nov 2026: the one unwritten axiom of the series, now written once.

\subsection{One complex field from two sheets}\label{sec:field}

Axiom A2 supplies a growth law $w=w_{0}e^{\kappa s}$ along the warped
direction and an involution $\iota$ with $\iota^{*}\kappa=-\kappa$
(M1 Definition~2.2). The two signs of $\kappa$ are the two \emph{sheets} of
the exponential, growth and decay, exchanged by $\iota$. M1 identifies
$\iota$ on the fibre algebra: at the drive-fixed corner $P_{1}=C_{0}$ it is
conjugation on $\Hquat$, whose commutative fixed subalgebra containing the
selected unit is $\mathbb{C}_{i}$ (M1 Propositions~5.4 and~5.6); at the three
cycled corners $P_{3}=C_{i}$, $P_{2}=C_{j}$, $P_{4}=C_{k}$ it is the inner
automorphism by the surviving unit, and the drive's three-cycle
$\Ad_{w}:i\to j\to k$ carries one corner's reading to the next
(M1 Proposition~5.6, certificate \texttt{L2\_Ad\_w\_cycles}).

Before the definition, three facts it rests on, each from M1, stated in
words. First, the third axiom says the fibre algebra --- the quaternions,
by M1 Proposition~5.2 --- acts on spinors (M1 Definition~2.5); so whatever
field the two sheets of the growth law make, it is a spinor-type object,
and its values live in the space on which the quaternions act. Second,
that space is fixed by the dimension: in twelve dimensions the Clifford
algebra is the algebra of $32\times32$ quaternionic matrices, by the
eightfold periodicity of Bott (twelve is four more than eight, and the
four-dimensional case is $\Hquat(2)$, so the twelve-dimensional case is
$\Hquat(2)$ enlarged sixteenfold); the space it acts on irreducibly is
therefore thirty-two quaternions, which is sixty-four complex numbers or
one hundred and twenty-eight real ones (\texttt{L7\_module\_dims}). Third,
the operation that exchanges the two sheets is the flow reversal $\tau$ of
Section~\ref{sec:handedness}: reverse the drive and turn the fibre by a
quarter. Done twice it is the fold of the torus that makes the pillowcase,
so on the pillowcase itself it is an involution.

\begin{definition}[the field $\psi$ and its shadow]\label{def:psi}
Let $M\cong\Hquat^{32}$ be the irreducible module of the Clifford algebra
of $M_{12}$. A \emph{sheet field} is a smooth section $\psi$ of $M$: at
every point of the shape, thirty-two quaternions
$\psi_{1},\dots,\psi_{32}$, varying smoothly from point to point. Let
$\hat\tau$ be the lift of the flow reversal $\tau$ to $M$ --- the operation
that carries a spinor at a point to a spinor at the reversed point, and
which, because $\tau$ reverses orientation, is an odd element of the Pin
group. The \emph{shadow} of $\psi$ is the sheet field
\begin{equation}\label{eq:equivariance}
  \bar\psi\;\equiv\;\hat\tau\psi\,,
\end{equation}
the field as seen after the exchange of growth and decay. Applying
$\hat\tau$ twice returns $\psi$ carried round the fold, which on the
pillowcase is $\psi$ itself up to the sign $\pm1$ that M1 Corollary~7.4
assigns to each corner. So on the pillowcase every sheet field splits into
an even part and an odd part under the exchange,
\[
  \psi\;=\;\tfrac12(\psi+\bar\psi)\;+\;\tfrac12(\psi-\bar\psi)\,,
\]
the first unchanged by exchanging growth and decay, the second reversed
by it. Finally, at a corner $C_{u}$ of the pillowcase the quaternions
collapse to the complex numbers $\mathbb{C}_{u}$ fixed by the corner
reading (M1 Propositions~5.4 and~5.6); the \emph{corner reading} of a
sheet field is its image there, thirty-two complex numbers instead of
thirty-two quaternions, and it is in that reading that the phase of
Section~\ref{sec:omega-shape} is defined.
\end{definition}
%% PI rulings 8 Nov 2026 (R-F1(Φ) adopted; Q17(b)): ψ is a section of the
%% Clifford module M ≅ ℍ³² (A3), the bar is the geometric lift τ̂ of the flow
%% reversal WITHOUT conjugation (so ω = ψ + τ̂ψ is quaternion-valued, carrying
%% the imaginary direction along the selected unit — R-τ11), and the
%% corner-algebra bundle 𝒞 is the corner reading of the module-valued field.

The two sheets of A2 are therefore one field: its $\hat\tau$-even part is
what growth and decay share, its $\hat\tau$-odd part is what distinguishes
them. That the field is module-valued rather than real or complex is not an
extra choice: M1 Proposition~7.3 reads the spinor data of A3 as the
$\Spinc$ module, and the selected imaginary unit of Proposition~7.2 is the
direction in $\mathrm{Im}\,\Hquat$ along which the even part carries its
imaginary component. In the corner reading each component is a complex
number $x+iy$ in $\mathbb{C}_{u}$ and $\hat\tau$ sends $y\mapsto-y$.
\rigD

\subsection{The even and odd variables, and the dilaton}\label{sec:omega}

\begin{definition}[$\omega$ and the odd variable]\label{def:omega}
\begin{equation}\label{eq:omega-def}
  \omega\;\equiv\;\psi+\bar\psi\,,
  \qquad
  \psi-\bar\psi\,,
\end{equation}
both sections of $M$. The first is the even part of the sheet field doubled,
the second the odd part doubled. $\omega$ is the Coherent Field of the
programme: it measures how much the two sheets agree. It is quaternion-valued,
and its imaginary part lies along the selected unit; since the odd variable
will enter the dilaton only through its modulus, a non-negative real number,
$\omega$ is the only one of the two that can give the dilaton a direction
(Section~\ref{sec:omega}). In the corner reading, where each component is a
complex number, $\omega=2\,\mathrm{Re}\,\psi$ and $\psi-\bar\psi=2i\,\mathrm{Im}\,\psi$.
\end{definition}

Two properties of $\omega$ are fixed at the level of the definition and are
the inputs to the next proposition. First, $\omega$ is the even
\emph{combination} of the sheets and not their modulus: when growth stands
exactly against decay, $\psi=-\bar\psi$, nothing in the geometry changes and
$\omega=0$. The modulus $|\psi|$ does not vanish there; the combination
does. Second, when the two sheets coincide, $\psi=\bar\psi$, there is no held
distinction between them, and we require that whatever couples the sheets to
the curvature vanish: this is the motivated condition that the coupling is a
coupling \emph{between} the sheets, not a coupling of either sheet alone.
\rigM
%% PI ruling R1 (8 Nov 2026): ω = ψ + ψ̄, with the reason as stated.
%% The second condition reproduces omega-definition-usage_v2.md §4 in-document.

\begin{proposition}[the lowest even invariant]\label{prop:parity}
Write $E=\tfrac12(\psi+\bar\psi)$ for the even part and $O=\tfrac12(\psi-\bar\psi)$
for the odd part of a sheet field. Consider scalars built by multiplying
powers of $E$ and powers of $\lvert O\rvert$ together, and ask for those
that are unchanged when growth and decay are exchanged, and that vanish
both when the sheets coincide ($O=0$) and when growth exactly opposes
decay ($E=0$). Every such scalar has total degree at least three, and the
one of degree three is, up to a real constant, $E\,\lvert O\rvert^{2}$.
In the programme's normalisation it is the dilaton
\begin{equation}\label{eq:PhiQC}
  \QCD\;=\;\omega\,\lvert\psi-\bar\psi\rvert^{2}\,,
\end{equation}
taken component by component in $M$: the $\alpha$-th component of $\QCD$
is $\omega_{\alpha}$ times the squared modulus of $(\psi-\bar\psi)_{\alpha}$.
\end{proposition}

\begin{proof}
We work in the corner reading, where each component is a complex number
$x+iy$ with $x=\mathrm{Re}\,\psi$ the even part and $y=\mathrm{Im}\,\psi$ the
odd part; the general case is the same argument in each component. A
monomial $x^{a}y^{b}$ changes sign under the exchange exactly when $b$ is
odd, because the exchange sends $y$ to $-y$ and leaves $x$ alone; so the
even monomials are those with $b$ even. The monomial vanishes identically
on the locus $y=0$ only if $y$ actually appears, that is $b\ge1$, and on the
locus $x=0$ only if $x$ appears, $a\ge1$. Put the three conditions together:
$b$ is even and at least one, so $b\ge2$; $a\ge1$; hence $a+b\ge3$, with
equality only when $a=1$ and $b=2$. The unique degree-three candidate is
therefore $xy^{2}$. In the original variables $x=\omega/2$ and
$y^{2}=\lvert\psi-\bar\psi\rvert^{2}/4$, so $xy^{2}=\tfrac18\,\omega\lvert\psi-\bar\psi\rvert^{2}$,
which is \eqref{eq:PhiQC} up to the constant $\tfrac18$. The classification
of exponents, the two vanishing criteria and the parity table of $\omega$,
$\psi-\bar\psi$, $\lvert\psi-\bar\psi\rvert^{2}$ and $\QCD$ are certified in
\texttt{M2\_L1\_parity.lean} (\texttt{L1\_unique}, \texttt{L1\_vanish\_Im},
\texttt{L1\_vanish\_Re}, \texttt{L1\_table\_*}). \rigR
\end{proof}

One remark on the module. The factor $\lvert\psi-\bar\psi\rvert^{2}$ is a
non-negative real number in each component, so it carries no direction in
the quaternions: whatever direction $\QCD$ has, it has because $\omega$ has
it. In particular $\QCD$ has an imaginary part only through $\omega$, and
$\omega$ is the unique carrier of the imaginary direction
(\texttt{L7\_imaginary\_square\_real}, which certifies the supporting fact
that the square of an imaginary quaternion is the negative of its squared
modulus). The real factor also commutes with everything, so no question of
the order of factors arises. The quaternionic structure that the fibre
paper of the series reads as the $Sp(1)$ flux on $\Stwo$ is carried by the
direction of $\omega$; this paper does not use it. \rigR

We record what the proposition does and does not decide. Parity alone does
not choose $\psi+\bar\psi$ over the modulus $\lvert\psi\rvert$ as the even
factor: both are unchanged by the exchange. The choice is the first property above, that $\omega$ vanishes
when growth stands against decay; with the modulus in its place $\QCD$ would
not vanish there. The proposition then shows that once the even variable is
$\omega$ and the two vanishing conditions are imposed, there is no further
freedom at lowest order.
%% Correction logged (handover): an earlier argument treating the fold and
%% conjugation as independent involutions was wrong; they are one involution.
%% PI ruling 9 Nov 2026: the square in Φ is the modulus squared |ψ−ψ̄|²
%% (the Source of Truth wrote (ψ−ψ̄)²; Corpus Errors entry 19).

The smallest case by hand: the drive-fixed corner $P_{1}$ is fixed by $\tau$
(Section~\ref{sec:handedness}), so $\bar\psi(P_{1})=\psi(P_{1})$ up to the
corner sign, so $\mathrm{Im}\,\psi(0)=0$ and
$\QCD(P_{1})=0$: the dilaton vanishes at the real corner, where there is one
sheet and nothing to hold apart.

\subsection{What is varied, and over what}\label{sec:scheme}

The field equation of this paper is obtained by variation and reduction. We
state the scheme once so that each later section can cite it.

\begin{definition}[the variational data]\label{def:scheme}
\begin{enumerate}[label=(\roman*),nosep]
\item \emph{The domain} is $M_{12}$ of M1: compact, connected, without
boundary (A1, M1 Definition~2.1).
\item \emph{The fields varied} are the twelve-dimensional metric $g$, with
the A2 warp $w=e^{A}$ depending on the base coordinates
$(\theta_{\psi},\theta_{\omega})$ (M1 Definition~2.2(i)); the dilaton
$\QCD$ of Section~\ref{sec:omega}, a section of the Clifford module $M$ on
$M_{12}$ with coupling $\zeta$ (Section~\ref{sec:equation} states how its kinetic term arises from the sheets);
the $\Spinc$ spinor of A3 on the bundle with determinant line
$L=\mathcal{O}(3)$ (M1 Proposition~7.3); and the $\Spinc$ connection.
\item \emph{The reduction} integrates over $\Keight$ in two stages: over the
six-dimensional fibre $\CPtwo\times\Stwo$, leaving a two-dimensional system
on the pillowcase base $\Ttwo/\Ztwo$ (Section~\ref{sec:warp}); and over the base, leaving
the four-dimensional equation on $S^{1}\times\Sigma$ (Section~\ref{sec:equation}). The field
equation of this paper is the twelve-dimensional one, on $M_{12}$; the
reduction is how its consequences on the base and on the external block
are read.
\item \emph{The scale} of $M_{12}$ is the declared input of M1
(Theorem~8.1, Proposition~8.2) and is not varied.
\end{enumerate}
\end{definition}
%% ARCH: section numbers are placeholders (§3 = WP3, §4 = WP4); renumber at
%% assembly. ROLLBACK 8 Nov 2026: the earlier sentence identifying Φ_QC with
%% the 4D image of the warp was the handover's error and is removed; Φ_QC is
%% a dilaton on M₁₂ (re-scoped R3.1).


%%=====================================================================
%% M2 — WP2 / R2.3: the two sectors are mirror images relative to the flow.
%% SURGICAL FRAGMENT (insertion only; no preamble). UK English. "We".
%% Architect draft, 8 Nov 2026, under PI rulings Q8 (Riemannian; time =
%% orbit parameter of the drive), Q9 (handedness relative to the flow,
%% derived in M2; Sol carries the arrow), Q10 (one warped direction, two
%% forced rates). Depends on Def. def:psi (WP0, revised) and M1 §4, §7.
%% Intended location: M2 Section 2, after the parity proposition.
%%=====================================================================

\subsection{The arrow, and why the shadow is not the image}\label{sec:handedness}

Nothing so far distinguishes $\psi$ from $\bar\psi$ except a sign under the
fold. If that were all, the second sector would be the first written
backwards, and $\lvert\psi-\bar\psi\rvert^{2}$ would measure a convention. This
subsection shows that it is not all: the operation relating the two sectors
reverses the orientation of the geometry, so that $\bar\psi$ is the mirror of
$\psi$ relative to the flow of A4, and the two carry opposite handedness. The
argument uses no sign in the metric. Under the fifth reading of
Section~\ref{sec:readings} the metric of $M_{12}$ is positive-definite and
signs no direction; the distinction between the sectors is made by the drive.

\paragraph{The Sol block.}
A4 places a flow on $T\times F$ whose return map is the drive
$M_{0}=\bigl(\begin{smallmatrix}2&1\\1&1\end{smallmatrix}\bigr)$
(M1 Theorem~4.2). The flow is therefore the suspension of $M_{0}$: the
three-manifold obtained from $\Ttwo\times[0,1]$ by identifying $(x,1)$ with
$(M_{0}x,0)$, with the flow $\partial_{s}$ along the interval. A torus bundle
over a circle with hyperbolic monodromy carries the Sol geometry
\cite{Thurston}, and in the eigenbasis of $M_{0}$ the Sol metric is
\begin{equation}\label{eq:sol-metric}
  \mathrm{d}\ell^{2}
  \;=\;e^{2\kappa s}\,\mathrm{d}x_{u}^{2}+e^{-2\kappa s}\,\mathrm{d}x_{s}^{2}
  +\mathrm{d}s^{2}\,,
  \qquad \kappa=\log\lambda=\log\varphi^{2}\,,
\end{equation}
where $x_{u},x_{s}$ are coordinates along the unstable and stable
eigenvectors $v_{u}=(\varphi,1)$ and $v_{s}=(-1,\varphi)$ of $M_{0}$, with
eigenvalues $\varphi^{2}$ and $\varphi^{-2}=2-\varphi$
(\texttt{L4\_eigen}). This is the factor $S^{1}_{\Sol}$ in the name of the
geometry, and it is A2's law seen in A4's drive: one direction grows as
$e^{\kappa s}$ and the other decays as $e^{-\kappa s}$, at the rate fixed by
the trace-three theorem. The two sheets of the exponential are the two
eigen-directions of the drive. \rigD

\paragraph{Reversing the flow.}
A symmetry of the Sol block that reverses the flow has the form
$(x,s)\mapsto(Ax,-s)$ with $A\in GL(2,\Z)$, and it is well defined on the
suspension iff $AM_{0}A^{-1}=M_{0}^{-1}$: reversing time inverts the
monodromy, and $A$ must carry the inverted monodromy back. Such an $A$
necessarily exchanges the two eigenlines, since it sends the eigenvalue
$\lambda$ to $\lambda^{-1}$. Two kinds exist.

\begin{proposition}[flow reversals]\label{prop:reversal}
Let $A\in GL(2,\Z)$ satisfy $AM_{0}A^{-1}=M_{0}^{-1}$.
\begin{enumerate}[label=(\roman*),nosep]
\item If $A^{2}=I$ then $\det A=-1$, and the reversal $(Ax,-s)$ preserves the
orientation of the Sol block.
\item If $\det A=+1$ then $A^{2}=-I$; the reversal has order four, its square
is the fold $x\mapsto-x$, and it reverses the orientation of the Sol block.
The quarter-turn $S=\bigl(\begin{smallmatrix}0&1\\-1&0\end{smallmatrix}\bigr)$
is such an $A$: $SM_{0}S^{-1}=M_{0}^{-1}$, $S^{2}=-I$, $Sv_{u}=-v_{s}$,
$Sv_{s}=v_{u}$.
\end{enumerate}
\end{proposition}

\begin{proof}
(i) Write the matrix as $A=\bigl(\begin{smallmatrix}a&b\\c&d\end{smallmatrix}\bigr)$
with integer entries. Squaring it and setting the result equal to the
identity gives four equations; the two off-diagonal ones read
$b(a+d)=0$ and $c(a+d)=0$. There are two cases. If $a+d\neq0$ then $b$ and
$c$ are both zero and $A$ is diagonal; but a diagonal matrix conjugating
$M_{0}$ to its inverse would have to send the top-left entry $2$ of $M_{0}$
to the top-left entry $1$ of $M_{0}^{-1}$, which forces $a=0$, contradicting
$a^{2}=1$ from the diagonal of $A^{2}=I$. So $a+d=0$, that is $d=-a$, and then
the determinant is $ad-bc=-a^{2}-bc=-(a^{2}+bc)=-1$, using the top-left
entry $a^{2}+bc=1$ of $A^{2}=I$. The reversal $(x,s)\mapsto(Ax,-s)$ then has
determinant $\det A\cdot(-1)=+1$: it preserves orientation
(\texttt{L4\_involution\_det}). (ii) Any $A$ that exchanges the two
eigenlines has zero diagonal in the eigenbasis, hence zero trace, and by
Cayley--Hamilton $A^{2}=(\operatorname{tr}A)A-(\det A)I=-I$ when
$\det A=1$. The quarter-turn $S$ is such an $A$: multiplying out,
$SM_{0}S^{-1}=M_{0}^{-1}$ and $S^{2}=-I$ (\texttt{L4\_S\_conj},
\texttt{L4\_S\_sq\_fold}); it sends $(\varphi,1)$ to $(1,-\varphi)$, which is
minus the stable eigenvector, and $(-1,\varphi)$ to $(\varphi,1)$
(\texttt{L4\_swap\_lines}); and the reversal built from it has determinant
$\det S\cdot(-1)=-1$, so it reverses orientation
(\texttt{L4\_tau\_reverses}). \rigR
\end{proof}
%% ARCH: (ii)'s trace-zero step is stated, not certified; the certificate
%% covers S and the general involutive case (i). If the PI wants (ii) in
%% full generality certified, it is a bounded enumeration — flag.

The smallest case by hand: $S$ sends $(\varphi,1)$ to $(1,-\varphi)=-(-1,\varphi)$,
the unstable line to the stable line. Running the flow backwards and turning
the fibre by a quarter exchanges growth and decay exactly.

\paragraph{The involution of A2 is the quarter-turn reversal.}
A2(iii) asks for an involution of the warped piece with $\iota^{*}\kappa=-\kappa$.
On the Sol block the only operations that send $\kappa$ to $-\kappa$ are the
flow reversals of Proposition~\ref{prop:reversal}, and among them only
case~(ii) reverses the orientation of the geometry. We take
\begin{equation}\label{eq:tau}
  \tau:(x,s)\longmapsto(Sx,-s)\,,\qquad \tau^{2}=\text{the fold},
\end{equation}
as the involution of A2. On the covering torus $\tau$ has order four; on the
pillowcase $\Ttwo/\{\pm1\}$, which is where M1 reads A2, $\tau^{2}$ is the
identity and $\tau$ is the involution M1 requires. M1's identification of
$\iota$ with the fold is therefore the statement about $\tau^{2}$; the present
paper refines it to $\tau$, whose square it is. \rigD
%% ARCH / REGISTER: this refines M1's identification (fold ↔ A2's ι) to
%% (τ ↔ A2's ι, τ² = fold). Consistent with M1 on the pillowcase; a genuine
%% refinement on the cover. Flagged for the Register and for M1 v3 wording.
%% Also consistent with the corpus's quarter-turn Spin^c holonomy e^{iπ/2}.

\paragraph{Handedness.}
A reflection of one direction of a Riemannian manifold lifts to the spinor
bundle as an element of the Pin group lying in the odd component, and every
odd element anticommutes with the chirality operator
$\Gamma=\Gamma_{1}\cdots\Gamma_{12}$ \cite{LawsonMichelsohn}. The lift
$\hat\tau$ of $\tau$ reverses orientation (Proposition~\ref{prop:reversal}),
so $\hat\tau\Gamma=-\Gamma\hat\tau$.

\begin{theorem}[the two sectors have opposite handedness]\label{thm:handedness}
Let $\psi$ be a sheet field of definite chirality, $\Gamma\psi=\epsilon\psi$,
and $\bar\psi\equiv\hat\tau\psi$ with $\hat\tau$ the Pin lift of $\tau$ acting on
the module $M$. Then $\Gamma\bar\psi=-\epsilon\bar\psi$.
Equivalently: the frame $(v_{u},v_{s},\partial_{s})$ in which $\psi$ grows
along the flow is carried by $\tau$ to a frame of the opposite orientation,
$(-v_{s},v_{u},-\partial_{s})$, in which $\bar\psi$ grows along the flow.
\end{theorem}

\begin{proof}
$\Gamma\bar\psi=\Gamma\hat\tau\psi=-\hat\tau\Gamma\psi=-\epsilon\hat\tau\psi
=-\epsilon\bar\psi$. The frame statement is \texttt{L4\_swap\_lines} together
with $s\mapsto-s$; the determinant of the change of frame is $-1$. \rigR
\end{proof}

This is the content of the word ``shadow''. The second sector is not the
first reflected in a mirror placed in space; it is the first reflected in the
flow. Reversing the drive exchanges growth with decay, and the geometry
records which is which: that record is the arrow of time, and the arrow is
what keeps $\psi$ and $\bar\psi$ apart. Had the flow no stable and unstable
direction, running it backwards would be indistinguishable from running it
forwards, there would be no arrow, and $\psi$ and $\bar\psi$ would be one
sector written twice. A4's positive entropy, which M1 used to force the
hyperbolic drive, is the same fact as the distinctness of the two sectors.

We record what has been used. The fifth reading (positive-definite metric)
makes the chirality operator a Riemannian one and the Pin argument available
in the form stated; no Lorentzian structure enters. The identification of
A2's involution with $\tau$ rather than with its square is the one
reading of this subsection, tagged \rigD\ above. The spinor lift of $\tau$
is a standard construction once $\tau$ is an isometry of the Sol metric,
which it is by \eqref{eq:sol-metric} since $S$ exchanges $x_{u}$ and
$x_{s}$ while $s\mapsto-s$ exchanges the exponents.

\paragraph{The warped direction.}
A2's warp and the Sol exponential of \eqref{eq:sol-metric} share the flow
direction $s$. They do not share a rate: the drive stretches the fibre torus
by $\varphi^{2}$ per circuit, a theorem of A4, while the warp of the
six-dimensional fibre over the base carries the exponent fixed in
Section~\ref{sec:warp}, a theorem of the flux. Two forced numbers are not an unforced
datum under M1's first minimality criterion, and the ratio of the two rates
is a pure number of the geometry with no input. Whether the base coordinate
$\theta_{\omega}$ of M1 is the flow coordinate $s$ is settled by the
reduction of Section~\ref{sec:equation} and is not assumed here. \rigM
%% PI ruling Q10. The ratio 24π/log φ² ≈ 78.4 is NOT stated in the text
%% (unchecked against the corpus); it stays in this comment until checked.

%%=====================================================================
%% M2 — WP3 / R3.1–R3.6 (RE-SCOPED, 8 Nov 2026): the field equation on M₁₂.
%% SURGICAL FRAGMENT (insertion only; no preamble). UK English. "We".
%% Under PI rulings Q1 (one gravitational term), Q8 (Riemannian; time =
%% flow parameter), Q12 (induced Spin^c connection, no gauge coupling), and
%% the re-scoped R3.1 (Φ_QC a dilaton on M₁₂ with coupling ζ; the warp is
%% the geometry's response, a = −ζΦ; Paper 2's 2D system is the base
%% restriction). Supersedes the withdrawn 4D-radion draft entirely (SOW §8).
%% Depends on WP0 (def:scheme, def:riemannian, def:omega, prop:parity) and
%% WP2 (sec:handedness). Placeholder: Section~\ref{sec:warp} = WP4 (warp exponent).
%%=====================================================================

\section{The field equation}\label{sec:equation}

The purpose of this section is to write the action on $M_{12}$ that the
axioms admit, to vary it, and to display the field equation that results.
In plain terms: we write down the one number that every configuration of
the shape and its fields is assigned, allow it to contain only what the
axioms force or the smallest option they leave open, and ask for the
configurations at which it does not change. The answer is Einstein's
equation with one additional term, and the section says which axiom is
responsible for each part: closure for the curvature and the vacuum term,
order for the matter and the force field, growth and decay for the
additional term. The equation is twelve-dimensional. Its restriction to
the base of $\Keight$ is the system in which Section~\ref{sec:warp} derives
the warp exponent, and its four-dimensional consequences are the subject
of later papers of the series; neither is needed to state it. We announce
the result.

\begin{theorem}[the Master Equation]\label{thm:master}
Under the action \eqref{eq:action} the metric of $M_{12}$ satisfies
\begin{equation}\label{eq:master}
\begin{gathered}
  G_{MN}\;=\;8\pi G_{\mathrm{eff}}\,T_{MN}\;+\;\zeta\,T^{\QCD}_{MN}\;+\;\Lambda_{\mathrm{eff}}\,g_{MN}\,,\\
  T^{\QCD}_{MN}\;\equiv\;\langle\nabla_{M}\QCD,\nabla_{N}\QCD\rangle_{M}-\tfrac12\,g_{MN}\,\langle\nabla\QCD,\nabla\QCD\rangle_{M}\,,
\end{gathered}
\end{equation}
where $\QCD=\omega\lvert\psi-\bar\psi\rvert^{2}$ is the dilaton of
Definition~\ref{def:omega}, a section of the Clifford module $M\cong\Hquat^{32}$,
$\langle u,v\rangle_{M}=\sum_{\alpha}\mathrm{Re}(u_{\alpha}\bar v_{\alpha})$ the
module inner product, $\zeta$ the coupling, $T_{MN}$ the conserved
stress tensor of the spinor sector of A3 together with the field strength
of the $\Spinc$ line, $8\pi G_{\mathrm{eff}}=\kappa_{12}^{2}$, and
$\Lambda_{\mathrm{eff}}$ collects the vacuum terms. The dilaton satisfies
\begin{equation}\label{eq:dilaton}
  \zeta\,\Box\QCD\;=\;V'(\QCD)\,,
\end{equation}
so that $\Box\QCD=0$ wherever the potential is stationary or absent.
\end{theorem}

In words, the tensor $T^{\QCD}_{MN}$ is built as follows. Take the rate of
change of the dilaton in direction $M$ and the rate of change in direction
$N$; these are each thirty-two quaternions, and their inner product over
the module, summing the real parts of the products component by component,
is a single real number $\langle\nabla_{M}\QCD,\nabla_{N}\QCD\rangle_{M}$. That is
the first term: the stress the dilaton's variation puts between directions
$M$ and $N$. The second term subtracts half the total squared rate of change
of the dilaton, in all directions at once, spread evenly over the metric
$g_{MN}$; it is the pressure the dilaton's gradient exerts isotropically. The
difference is the standard stress tensor of a scalar field, familiar from
every textbook, with the one change that the squares are inner products
over the module. For a single real component it is exactly
$\nabla_{M}\QCD\nabla_{N}\QCD-\tfrac12g_{MN}\lvert\nabla\QCD\rvert^{2}$.

Read term by term, \eqref{eq:master} is Einstein's equation, $G_{MN}=8\pi G\,T_{MN}+\Lambda g_{MN}$,
with the one additional term $\zeta T^{\QCD}_{MN}$. The curvature on the
left and the vacuum term exist because the shape closes on itself (A1);
the matter and force term $T_{MN}$ exists because order matters (A3) and
the algebra of order acts on spinors; the additional term exists because
the shape grows and decays as a function of itself (A2), its field being
written on the two sheets of that law; and the two sheets are two, rather
than one sheet and its reflection, because the shape never stands still
(A4, Section~\ref{sec:handedness}). Einstein's equation is the case in
which the second and fourth axioms are absent.

Two readings are fixed now. By Definition~\ref{def:riemannian} the metric
is positive-definite and every ``time'' statement below is a statement
along the flow. And $\zeta$ in \eqref{eq:master} is a coupling constant of
the twelve-dimensional theory; Section~\ref{sec:warp} derives its value, $12\pi$, as the
exponent by which the warp of A2 falls across a half-period of the base
when the geometry responds to the dilaton's profile. The number is one
number seen in two places, and the equation displays it once.
%% Re-scoped R3.1. The withdrawn 4D radion computation is not this object.

\subsection{The action}\label{sec:action}

\begin{definition}[the action on $M_{12}$]\label{def:action}
\begin{equation}\label{eq:action}
\begin{aligned}
  S\;=\;\frac{1}{2\kappa_{12}^{2}}\int_{M_{12}}\!\sqrt{g}\,
  \Bigl[\,R\;-\;2\Lambda_{12}\;-\;\tfrac14\,F_{MN}F^{MN}
  \;-\;&\zeta\,\langle\nabla\QCD,\nabla\QCD\rangle_{M}\;-\;2V(\QCD)\Bigr]\\
  &+\;S_{\mathrm{spinor}}[g,\Psi,A]\,,
\end{aligned}
\end{equation}
with $F=\mathrm{d}A$ the curvature of the induced connection on the
determinant line $L=K^{-1}$ (Section~\ref{sec:connection}), $\QCD$ the
dilaton of Definition~\ref{def:omega}, $V$ its potential, and
$S_{\mathrm{spinor}}$ the Dirac action of the $\Spinc$ spinor of A3 on the
bundle of M1 Proposition~7.3.
\end{definition}
%% CONVENTION RULED (PI, 9 Nov 2026, M2-7): the Φ kinetic term sits inside
%% the 1/2κ² bracket with coefficient ζ; this reproduces the canonical display
%% with no κ² on T^Φ. CFM v12 l.1525's −½ζ outside the bracket is the
%% superseded normalisation (Corpus Errors entry 20).

In words: the number $S$ assigned to a configuration is the integral over
the whole twelve-dimensional shape of five contributions, all weighted by
the volume element $\sqrt{g}$ and divided by twice the gravitational
constant $\kappa_{12}^{2}$. The first, $R$, is the scalar curvature --- how
much the shape bends at each point. The second, $-2\Lambda_{12}$, is the
energy of empty volume, a constant per unit volume. The third,
$-\tfrac14F_{MN}F^{MN}$, is the energy of the field strength of the
$\Spinc$ line, the one abelian field the geometry carries. The fourth,
$-\zeta\langle\nabla\QCD,\nabla\QCD\rangle_{M}$, is the energy stored in the
dilaton's variation from point to point, the module inner product of its
gradient with itself, multiplied by the coupling $\zeta$. The fifth,
$-2V(\QCD)$, is the potential energy of the dilaton's value. Added to all
of these is the action of the spinor field of the third axiom, the Dirac
action $S_{\mathrm{spinor}}$, which depends on the metric, the spinor
$\Psi$ and the connection $A$. Each term is accounted for by an axiom or a
reading, and we go through them.

\paragraph{One gravitational term.}
By A1 the domain is compact without boundary (M1 Definition~2.1), so the
action has no boundary terms and its stationary points are Euler--Lagrange
equations with no boundary data \cite{Wald}. The divergence-free symmetric
tensors built from the metric and two derivatives in twelve dimensions
form the five-term Lovelock family \cite{Lovelock1971}; every term beyond
the first carries an independent dimensionful coefficient, an unforced
datum under criterion~(1) of M1 Definition~3.2, excluded on the ground that
excluded the twisted $\Spinc$ lines in M1 Proposition~7.3. The Einstein
term survives because without it there is no equation; its coefficient is
the declared scale of M1 Theorem~8.1. Once one term linear in second
derivatives is admitted its uniqueness holds in every dimension
\cite{Vermeil1917,Weyl1922,Cartan1922}. \rigD
%% PI ruling Q1.

\paragraph{The connection of the line.}\label{sec:connection}
M1 Proposition~7.3 selects the $\Spinc$ structure induced by the
almost-complex structure of A3's selected unit, with determinant line
$L=K^{-1}$ (M1 Definition~2.5). The connection on $L$ is part of the
geometry's own data, and its field strength is the one abelian field
strength the geometry carries. Its energy therefore enters
\eqref{eq:action} inside the gravitational bracket, weighted by
$1/2\kappa_{12}^{2}$: a separate gauge coupling for a connection the
geometry already owns would be an unforced datum, excluded by
criterion~(1). What is \emph{not} fixed by the metric is the flux: the
quanta of $F$ through the cycles of $\Keight$ are topological data, with
total $N=\chi(\Keight)$ by the tadpole condition (Section~\ref{sec:warp}), and they sit
on the base cycle because it is the only cycle with cone points at which a
localised flux can satisfy the first-order equation of Section~\ref{sec:warp}. \rigD
%% PI rulings Q12/Q14 (weak (c), 8 Nov 2026). Consistent with CFM v12 l.1419
%% (Maxwell term inside the bracket) and Paper I Addendum §4.4–4.6 (κ²_{2,eff}
%% = κ₈²/(Vol₀ e^{6A₀}); twelve units on the base cycle). The identification
%% F = Ricci form, briefly adopted, is withdrawn (Minimality Register M2-2).

\paragraph{The dilaton term, and why the sheets enter through $\QCD$ alone.}
The fundamental fields are the sheets $\psi,\bar\psi$ and the Coherent
Field $\omega$ (Section~\ref{sec:variables}); the spinor's own dynamics are
the Dirac action in $S_{\mathrm{spinor}}$. The dilaton
$\QCD=\omega\lvert\psi-\bar\psi\rvert^{2}$ is composite and module-valued, so a
two-derivative term in it is a particular interaction among the
fundamental fields and their gradients, and writing it as the only such
term is a statement: the geometry couples to the sheets through $\QCD$ and
through nothing else. Which quadratic form in $\nabla\QCD$ appears is not a
further choice: the invariant symmetric forms on an irreducible
quaternionic-unitary module form a one-dimensional space, so covariance
under the fibre algebra forces the module inner product
$\langle\nabla_{M}\QCD,\nabla_{N}\QCD\rangle_{M}$ --- the trace of the square,
not the square of the trace. For a vacuum symmetric under the module the
trace is $32$ times a single component, and the three readings of
``$\zeta\nabla\QCD\nabla\QCD$'' for a module-valued field --- one component,
the trace, the square of the trace --- carry the weights $1$, $32$ and $1024$
(\texttt{L7\_three\_readings}); only the second is covariant. \rigR
We adopt this by the standing principle of Section~\ref{sec:readings}, and
the reason it is the minimal coupling is Proposition~\ref{prop:parity}. If curvature is to respond to the sheets
through a scalar that is even under the fold and vanishes both where the
sheets coincide and where growth stands against decay, $\QCD$ is the unique
lowest-order candidate, and $\lvert\nabla\QCD\rvert^{2}$ is then the unique
two-derivative term with those properties; the competing terms fail the
vanishing conditions ($\lvert\nabla\omega\rvert^{2}$ does not vanish at
$\psi=\bar\psi$; $\lvert\nabla\psi\rvert^{2}$ vanishes nowhere). This is
A2 read on the metric: the geometry grows and decays in response to the
asymmetry of the sheets weighted by their coherence, and to nothing else.
$\zeta$ is the strength of that one coupling and is never absorbed into a
redefinition of the field. \rigD\ (minimality)
%% The corpus posits S_Φ (CFM l.1525; Paper 2 §7.3 "microscopic origin a
%% separate question"). M2 names the step and gives the parity theorem as
%% its reason. Tagged [D]: minimal, not the only conceivable, coupling.

\paragraph{The potential and the vacuum term.}
$\Lambda_{12}$ is the energy of empty twelve-volume; $V(\QCD)$ is the
dilaton's potential, whose form on the base is obtained in Section~\ref{sec:warp} from
the fibre curvature and $\Lambda_{12}$ when the action is evaluated on the
warped configuration. Neither is a free parameter of this paper:
$\Lambda_{12}$ is fixed by the requirement that the base system of
Section~\ref{sec:warp} admit a solution with the given topology, and $V$ is derived
there. Earlier work in the programme took $\Lambda_{12}=0$ as an axiom of
the construction; in this series that statement is the case in which the
solvability condition returns zero, to be established or refuted where the
fibre radii are derived, and it is not assumed here. \rigD
%% PI ruling (Λ₁₂ by stationarity, 8 Nov 2026).

\subsection{Variation}\label{sec:variation}

\begin{proof}[Proof of Theorem~\ref{thm:master}]
Varying \eqref{eq:action} with respect to $g^{MN}$,
\begin{equation}\label{eq:var}
\begin{aligned}
  G_{MN}+\Lambda_{12}g_{MN}
  \;=\;&\zeta\Bigl(\langle\nabla_{M}\QCD,\nabla_{N}\QCD\rangle_{M}-\tfrac12 g_{MN}\langle\nabla\QCD,\nabla\QCD\rangle_{M}\Bigr)
  \;-\;V(\QCD)\,g_{MN}\\
  &+\;\tfrac12\Bigl(F_{MP}F_{N}{}^{P}-\tfrac14 g_{MN}F^{2}\Bigr)
  \;+\;\kappa_{12}^{2}\,T^{\mathrm{sp}}_{MN}\,,
\end{aligned}
\end{equation}
where $T^{\mathrm{sp}}_{MN}=-\frac{2}{\sqrt g}\frac{\delta S_{\mathrm{spinor}}}{\delta g^{MN}}$.
The first bracket is $T^{\QCD}_{MN}$ with coefficient $\zeta$ and nothing
else; it is the canonical stress tensor of a module-valued scalar, the
trace over $M$ of the scalar form, and its form is forced by the kinetic
term, not chosen. For a single component it reduces to
$\nabla_{M}\QCD\nabla_{N}\QCD-\tfrac12g_{MN}\lvert\nabla\QCD\rvert^{2}$, which is
the display the programme has used; every value of $\langle\QCD\rangle$ in
later papers is a trace over the thirty-two components. Writing
$8\pi G_{\mathrm{eff}}T_{MN}\equiv\kappa_{12}^{2}T^{\mathrm{sp}}_{MN}+\tfrac12(F_{MP}F_{N}{}^{P}-\tfrac14g_{MN}F^{2})$
and $\Lambda_{\mathrm{eff}}g_{MN}\equiv-\bigl(\Lambda_{12}+V(\QCD)\bigr)g_{MN}$
gives \eqref{eq:master}. Varying with respect to $\QCD$ gives
\eqref{eq:dilaton}. \rigR
\end{proof}

The identification of $\Lambda_{\mathrm{eff}}$ carries the dilaton's
potential, so it is constant only where $V$ is stationary; this is the
reading under which the canonical display and \eqref{eq:dilaton} are
consistent, as the next proposition shows.

\begin{proposition}[consistency]\label{prop:bianchi}
$\nabla^{M}T^{\QCD}_{MN}=\langle\Box\QCD,\nabla_{N}\QCD\rangle_{M}$. Hence the contracted
Bianchi identity, the conservation of $T_{MN}$ and \eqref{eq:master} imply
\[
\zeta\langle\Box\QCD,\nabla_{N}\QCD\rangle_{M}=\nabla_{N}V=\langle V'(\QCD),\nabla_{N}\QCD\rangle_{M}\,,
\]
which is \eqref{eq:dilaton} on the support of $\nabla\QCD$; and where
$V'=0$, $\Box\QCD=0$ there. The computation is componentwise and the trace
is linear.
\end{proposition}

\begin{proof}
Take the divergence of the stress tensor, that is, contract its first
index with a derivative. The first term, the inner product of two
gradients, gives by the product rule two pieces: the Laplacian of the
dilaton paired with its gradient, $\langle\Box\QCD,\nabla_{N}\QCD\rangle_{M}$,
and the gradient paired with the second derivative,
$\langle\nabla^{M}\QCD,\nabla_{M}\nabla_{N}\QCD\rangle_{M}$. The second term,
half the squared gradient spread over the metric, gives on differentiation
$\langle\nabla^{M}\QCD,\nabla_{N}\nabla_{M}\QCD\rangle_{M}$, the gradient
paired with the second derivative in the other order. For a scalar the two
orders of covariant differentiation agree, so those two pieces cancel and
only $\langle\Box\QCD,\nabla_{N}\QCD\rangle_{M}$ survives:
\begin{align*}
\nabla^{M}\bigl(\langle\nabla_{M}\QCD,\nabla_{N}\QCD\rangle_{M}\bigr)-\tfrac12\nabla_{N}\langle\nabla\QCD,\nabla\QCD\rangle_{M}
&=\langle\Box\QCD,\nabla_{N}\QCD\rangle_{M}\,.
\end{align*}
Now take the divergence of both sides of \eqref{eq:master}. The left side
vanishes by the contracted Bianchi identity, the geometric fact that the
Einstein tensor is automatically divergence-free. The matter term vanishes
because $T_{MN}$ is conserved (Section~\ref{sec:matter}). The vacuum term
contributes the gradient of $\Lambda_{\mathrm{eff}}$, which is the gradient
of $-V(\QCD)$, that is $-\langle V'(\QCD),\nabla_{N}\QCD\rangle_{M}$. What is
left is
\[
\zeta\langle\Box\QCD,\nabla_{N}\QCD\rangle_{M}=\langle V'(\QCD),\nabla_{N}\QCD\rangle_{M}\,,
\]
which is \eqref{eq:dilaton} wherever the dilaton is actually varying; and
where the potential is flat, $V'=0$, it says the dilaton satisfies the wave
equation $\Box\QCD=0$ there. The computation is componentwise and the trace
is linear, so nothing changes in passing from one component to the module.
\rigR
\end{proof}
%% The F1 note's "□Φ = 0 on supp ∇Φ" is the V' = 0 case; flag for the note.

\begin{proposition}[the dilaton stress tensor along the flow]\label{prop:stiff}
Let $\QCD$ depend only on the flow parameter $s$ of A4. In an orthonormal
frame $(\partial_{s},e_{1},\dots,e_{11})$ the tensor $\zeta T^{\QCD}$ is
diagonal with
\begin{equation}\label{eq:stiff}
  \zeta T^{\QCD}_{ss}=\tfrac12\zeta\langle\dot\QCD,\dot\QCD\rangle_{M}\,,\qquad
  \zeta T^{\QCD}_{ii}=-\tfrac12\zeta\langle\dot\QCD,\dot\QCD\rangle_{M}\quad(i=1,\dots,11)\,.
\end{equation}
\end{proposition}

\begin{proof}
If the dilaton depends only on the flow parameter, its gradient points
along the flow and nowhere else: $\nabla\QCD=\dot\QCD\,\partial_{s}$, where
the dot is the derivative along the flow. Its squared length is then just
$\langle\dot\QCD,\dot\QCD\rangle_{M}$. In the component along the flow the
first term of the stress tensor contributes the whole of that, and the
second term subtracts half of it, leaving a half. In any direction
transverse to the flow the first term contributes nothing, because the
gradient has no component there, and the second term subtracts half the
squared length, leaving minus a half. Off-diagonal components vanish for
the same reason (\texttt{L2\_flow\_components}). \rigR
\end{proof}

The reading. In a Lorentzian presentation, with the flow parameter
continued to a signed time, the two entries of \eqref{eq:stiff} are an
energy density and an isotropic pressure with
$\rho_{\QCD}=p_{\QCD}=\tfrac12\zeta\langle\dot\QCD,\dot\QCD\rangle_{M}$, the stiff equation of
state, and \eqref{eq:dilaton} with $V'=0$ gives $\rho_{\QCD}\propto a^{-6}$
on an isotropic background. These are the statements the programme uses.
Under Definition~\ref{def:riemannian} they are the components
\eqref{eq:stiff} read along the flow, the transverse sign carried by the
continuation; the equation is \eqref{eq:master} in either reading. \rigD

\subsection{The base restriction}\label{sec:base-restriction}

Section~\ref{sec:warp} works on the base $\Ttwo/\Ztwo$ with the fibre warped. We record
here, so that Section~\ref{sec:warp} may cite it, what the Master Equation becomes
there. Evaluate \eqref{eq:action} on the configuration of
Definition~\ref{def:scheme}: metric $g_{4}+e^{2a(\theta)}\hat g_{6}+g_{2}(\theta)$,
dilaton $\QCD=\QCD(\theta)$, the connection induced. The fibre volume
contributes $e^{6a}$, its curvature $e^{-2a}\hat R_{6}$, and the
warped-product scalar curvature contributes the gradient terms of $a$
\cite[Ch.~7]{ONeill1983}; the result is a two-dimensional action on the
base of the form
\begin{equation}\label{eq:2D}
  S_{2}\;=\;\int_{\Ttwo/\Ztwo}\!\sqrt{g_{2}}\,
  \Bigl[\frac{1}{2\kappa_{2}^{2}}R_{2}\;-\;\tfrac12\lvert F\rvert^{2}
  \;-\;\mathcal{V}(a,\QCD)\;-\;\mathcal{K}(a,\nabla a,\nabla\QCD)\Bigr]\,,
\end{equation}
with $\kappa_{2}^{-2}$ the fibre-volume-weighted gravitational
coefficient, $\mathcal{V}$ the potential assembled from $\Lambda_{12}$,
$\hat R_{6}$ and $V(\QCD)$ with the weights $e^{6a}$ and $e^{4a}$, and
$\mathcal{K}$ the gradient terms of $a$ and $\QCD$. The warp of A2 is the
geometry's response to the dilaton: on the solutions of Section~\ref{sec:warp} the two
profiles are tied by
\begin{equation}\label{eq:a-Phi}
  a\;=\;-\zeta\,\QCD\,,
\end{equation}
with $\QCD$ falling by one unit across a half-period, so that the warp
exponent across the half-period is $\zeta$ and \eqref{eq:2D} collapses to
the single-modulus Einstein--Maxwell system with potential that Section~\ref{sec:warp}
solves. That system is therefore the base restriction of
\eqref{eq:master}, not an axiom. \rigD
%% Discharges Paper 2's Σ-2 as a restriction of the Master Equation. The
%% relation eq:a-Phi is Paper 2 §7.3's definition of Φ as the modulus of the
%% warped direction, here read as the solution relation; its derivation as a
%% property of the solutions, with the boundary configuration Φ = 1, is in
%% Section 4 (WP4), which reproduces Paper 2 §4 on this footing.

\subsection{The matter sector}\label{sec:matter}

$T_{MN}$ in \eqref{eq:master} is the stress tensor of the $\Spinc$ Dirac
field of A3 on the bundle with $L=\mathcal{O}(3)$ (M1 Proposition~7.3) with
the Clifford action of the fibre algebra (M1 Definition~2.5), together
with the field strength of the induced connection. The one property of
$T_{MN}$ used in this paper is its conservation, which follows from the
diffeomorphism invariance of $S_{\mathrm{spinor}}$ on the closed manifold
$M_{12}$ (A1) by the standard argument \cite{Wald}; for the field strength
it is the Bianchi identity $\mathrm{d}F=0$ and the Maxwell equation. The
content of $T_{MN}$, its decomposition into generations and its masses,
belongs to later papers of the series and is not opened here. \rigD

\subsection{What is derived and what is read}\label{sec:summary3}

The form of \eqref{eq:master} is a theorem of the action \eqref{eq:action}
(\rigR). The action is fixed by one derivation and three readings: the
single gravitational term (minimality, \rigD), the induced connection
(\rigD), the coupling of the geometry to the sheets through $\QCD$ alone
(\rigD, with Proposition~\ref{prop:parity} as its reason), and the
Riemannian metric with time along the flow (Definition~\ref{def:riemannian}).
The coupling $\zeta$ is a constant of the theory whose value is derived in
Section~\ref{sec:warp}; the relation \eqref{eq:a-Phi} between the warp and the dilaton
is a property of the solutions established there.

%%=====================================================================
%% M2 — WP4: the warp exponent and the value of ζ. SURGICAL FRAGMENT.
%% UK English. "We". Reproduces, in-document and without citation, the
%% derivation of Paper 2 §§3–4 and Paper I Addendum §§3–4 (fossil), on the
%% footing of M2 §3 (the base system is the restriction of the Master
%% Equation). PI rulings: Q12/Q14 (weak (c)), Q15 (direction: τ-symmetric,
%% j/k convention; τ₂ physical), Q16 (jump exact, constant slope a boundary
%% condition), Λ₁₂ by stationarity (M2-6). Placeholders: §3 = WP3, §2 = WP0/WP2.
%%=====================================================================

\section{The warp exponent and the value of $\zeta$}\label{sec:warp}

The purpose of this section is to solve the base restriction of the Master
Equation, \eqref{eq:2D}, for the profile of the warp across a half-period
of the pillowcase, and to show that the exponent by which the warp falls is
$N\pi$ with $N$ the number of flux units, that $N=\chi(\Keight)=12$, and
hence that the coupling of Section~\ref{sec:equation} has the value
$\zeta=12\pi$. We announce the two results.

\begin{theorem}[exact warp exponent per flux unit]\label{thm:warp-exponent}
For the base system \eqref{eq:2D} on $\Ttwo/\Ztwo$ with $N$ units of
quantised flux localised at a cone point, the warp exponent across the
half-period transverse to the wall on which the flux sits satisfies
exactly
\begin{equation}\label{eq:kL-Npi}
  kL\;=\;N\pi\,,
\end{equation}
where $k$ is the constant decay rate of the warp exponent between the
walls and $L$ the corresponding half-period. The factor $\pi$ is the
deficit angle at the cone point.
\end{theorem}

\begin{corollary}[the value of $\zeta$]\label{cor:zeta}
With $N=\chi(\Keight)=12$ (Proposition~\ref{prop:flux}) and the dilaton
normalised to fall by one unit across the half-period (\eqref{eq:a-Phi}),
\begin{equation}\label{eq:zeta}
  \zeta\;=\;kL\;=\;\chi(\Keight)\,\pi\;=\;12\pi\,.
\end{equation}
\end{corollary}

\subsection{The topological integers}\label{sec:integers}

\begin{proposition}\label{prop:chi-pillowcase}
$\chi(\Ttwo/\Ztwo)=2$, and the total curvature of the pillowcase is $4\pi$,
concentrated as a deficit of exactly $\pi$ at each of the four cone points.
\end{proposition}

\begin{proof}
We give two independent routes to the Euler characteristic, and then the
curvature. First route, by the cone points. The Euler characteristic of a
torus is zero. Folding it in half by the map $z\mapsto-z$ would halve that
to zero again, except at the four points the fold leaves fixed, where the
quotient is not smooth: each fixed point contributes a correction of
$1-\tfrac12=\tfrac12$, because a neighbourhood of it is a cone of angle
$\pi$, half a full turn. Four corrections of a half make two, so
$\chi=0+4\cdot\tfrac12=2$ (M1 Theorem~5.7, \texttt{L2\_chi\_pillowcase}).
Second route, by counting holes. The folded torus is, as a topological
space, a sphere --- stretch the pillowcase out and it is a ball with four
puckered points --- and a sphere has one connected piece, no independent
loops and one enclosed cavity, Betti numbers $(1,0,1)$, so $\chi=1-0+1=2$.
The two routes agree. For the curvature, the Gauss--Bonnet theorem for
orbifolds \cite{Kawasaki1979} says the total curvature is $2\pi$ times the
Euler characteristic, $4\pi$. The pillowcase is flat everywhere except at
the four cone points, so all $4\pi$ of curvature sits at them, $\pi$ at
each --- which is exactly the angle missing at a cone of angle $\pi$
(\texttt{L5\_deficit}). \rigR
\end{proof}

\begin{proposition}\label{prop:chi-K8}
$\chi(\Keight)=\chi(\CPtwo)\,\chi(\Stwo)\,\chi(\Ttwo/\Ztwo)=3\cdot2\cdot2=12$.
\end{proposition}

\begin{proof}
Multiplicativity of $\chi$ over fibrations with compact fibre; the warp
does not change the additive structure of cohomology
(M1 Theorem~7.6, \texttt{L3\_chi\_K8}). \rigR
\end{proof}

\begin{proposition}[the flux number]\label{prop:flux}
Tadpole cancellation on $\Keight$ fixes the number of quantised flux units
of the determinant line to $N_{\mathrm{flux}}=\chi(\Keight)=12$, and all
twelve sit on the base cycle.
\end{proposition}

\begin{proof}
On a compact space the total charge sourced must cancel. The gravitational
contribution to the tadpole is the Euler density, whose integral is
$\chi(\Keight)$; the flux contribution is integrally quantised and must
balance it, so $N_{\mathrm{flux}}=\chi(\Keight)=12$. Cross-check from the
parity structure: the individual quanta are odd on the non-spin factor
$\CPtwo$ and even on the spin factor $\Stwo$, by $c_{1}\equiv w_{2}\bmod2$
(M1 Proposition~7.3, \texttt{L4\_flux\_parity}), and the values
$(Q_{c},Q_{w})=(3,6)$ of M1 Proposition~8.2 are the ones the constraint
admits. The twelve units sit on the base because a localised flux
satisfying the first-order equation of Section~\ref{sec:bps} needs a cone
point to sit at, and the base is the only cycle that has them. The
smallest case: one unit at one corner balanced by one unit of integrated
Euler density on its cone neighbourhood; the general case is the sum of
twelve such balances. \rigD
\end{proof}
%% Weak (c), Q14: flux quanta are topological data; the placement argument is
%% the Addendum's §4.6 reason, reproduced. Minimality Register M2-2.

\subsection{The base system and its single modulus}\label{sec:base-system}

Section~\ref{sec:base-restriction} obtained the base restriction
\eqref{eq:2D} of the Master Equation: two-dimensional gravity with
coefficient $\kappa_{2}^{-2}$, the energy $\tfrac12\lvert F\rvert^{2}$ of
the line's field strength, a potential, and the gradient terms of the
fibre warp $a$ and the dilaton $\QCD$. Two facts reduce it to a system with
a single modulus.

\paragraph{The fibre decouples from the jump.}
$\CPtwo$ with the Fubini--Study metric and $\Stwo$ with the round metric
are Einstein, $\mathrm{Ric}=(6/R_{c}^{2})g$ and $\mathrm{Ric}=(1/R_{w}^{2})g$
(M1 Propositions~7.1, 7.2). For a warp depending only on the base
coordinates the warped-product curvature \cite{BishopONeill1969} therefore
contributes to the base equation only smooth scalar terms in $a$ --- a
potential and a modification of the gradient term --- with no mixed
base--fibre components. Smooth terms cannot change a jump condition at a
cone point; they change the shape of the profile between cone points.
The effective two-dimensional gravitational coefficient is the
fibre-volume-weighted one,
\begin{equation}\label{eq:kappa-eff}
  \kappa_{2}^{2}\;=\;\frac{\kappa_{12}^{2}}{V_{4}\,\mathrm{Vol}_{0}\,e^{6a_{0}}}\,,
  \qquad \mathrm{Vol}_{0}=\mathrm{Vol}(\CPtwo)\,\mathrm{Vol}(\Stwo)\,,
\end{equation}
with $V_{4}$ the external volume and $a_{0}$ the warp at the wall. \rigD

\paragraph{The dilaton and the warp are one modulus.}
On the solutions constructed below the dilaton's profile and the fibre
warp are tied by \eqref{eq:a-Phi}, $a=-\zeta\QCD$, with $\QCD$ the
canonically normalised modulus of the warped direction and $\QCD=1$ the
configuration at the far wall. This is the definition of the
normalisation of $\QCD$ on the base, and it is what makes the coupling
$\zeta$ of \eqref{eq:master} and the exponent of \eqref{eq:kL-Npi} one
number: a profile of $\QCD$ that falls by one unit across the half-period
is a warp exponent that falls by $\zeta$. In the fossil form of this
construction the transverse base length itself was written with the
factor $e^{2\zeta\QCD}$ while the fibre carried $e^{2a}$; under
\eqref{eq:a-Phi} these are two expressions of the same exponent and not
two warps. \rigD
%% Q16 reconciliation: Paper 2 §7.3 (a = −ζΦ, Φ = 1 boundary) adopted as the
%% normalisation; Addendum eq. 8D-metric's e^{2ζΦ}dθ_ω² read as the same
%% exponent on the transverse length. Flag: "canonically normalised" here
%% means the normalisation that makes ΔΦ = 1 per half-period, fixed by the
%% boundary configuration; M2-7 (bracket convention) remains with the PI.

With these two facts \eqref{eq:2D} is the single-modulus system
\begin{equation}\label{eq:2D-single}
  S_{2}\;=\;\int_{\Ttwo/\Ztwo}\!\sqrt{g_{2}}\,
  \Bigl[\frac{1}{2\kappa_{2}^{2}}R_{2}\;-\;\tfrac12\lvert F\rvert^{2}\;-\;V(a)\Bigr]\,,
\end{equation}
with $V(a)$ the bulk potential assembled from $\Lambda_{12}$, the fibre
curvature and the dilaton potential, continuous on the base. It is the
restriction of the Master Equation and not an axiom. \rigD

\subsection{Bogomol'nyi completion}\label{sec:bps}

Let $\theta$ be the base coordinate transverse to the wall on which the
flux sits, and consider a static profile $a(\theta)$. The energy
functional of \eqref{eq:2D-single} admits a completion of the square once
$V(a)$ balances the geometric bulk terms,
\begin{equation}\label{eq:square}
  \mathcal{E}\;\propto\;\int\sqrt{g_{2}}\,
  \Bigl(\partial_{\theta}a\;\pm\;\frac{\kappa_{2}^{2}}{2}\,\frac{\lvert F_{\psi\omega}\rvert}{\sqrt{g_{2}}}\Bigr)^{2}
  \;\ge\;0\,,
\end{equation}
saturated when the square vanishes. The balance of $V$ against the bulk
terms is the familiar one of warped five-dimensional models
\cite{RS1999}; here it is not a tuning but the condition that
$\Lambda_{12}$ be the value at which the base system admits a solution
with the given topology, which is how $\Lambda_{12}$ is fixed in this
series (Section~\ref{sec:action}). The negative root, a decaying warp,
gives the first-order equation
\begin{equation}\label{eq:bps}
  \partial_{\theta}a\;=\;-\frac{\kappa_{2}^{2}}{2}\,\frac{\lvert F_{\psi\omega}\rvert}{\sqrt{g_{2}}}\,,
\end{equation}
and everything below is an exact consequence of it. \rigR\ (as a theorem of
\eqref{eq:2D-single}); the balance condition is \rigD.
%% Minimality Register M2-6 (Λ₁₂ by stationarity) is where the "tuning" lives.

\subsection{Flux localisation and quantisation}\label{sec:flux}

Dirac quantisation over the orbifold is $\frac{1}{2\pi}\int F=N$. For $N$
units at the drive-fixed corner $P_{1}=(0,0)$ --- the one corner fixed by
both the fold and the drive (M1 Propositions~4.3, 5.6) --- the density is
\begin{equation}\label{eq:flux-density}
  F_{\psi\omega}\;=\;N\pi\,\frac{\delta(\theta_{\psi})\,\delta(\theta_{\omega})}{\sqrt{g_{2}}}\,,
\end{equation}
checked directly: over the fundamental domain
$\frac{1}{2\pi}\int F=N/2$, and the domain is half of the covering torus,
so the total over the cover is $2\cdot N/2=N$ (\texttt{L5\_FD\_check}).
\rigR

\subsection{Solving the first-order equation}\label{sec:solve}

Away from the corner the localised flux vanishes and the flow is driven by
the constant bulk potential, giving the linear profile
\begin{equation}\label{eq:bulk}
  \partial_{\theta}a=-k\qquad\Longrightarrow\qquad a(\theta)=a(0)-k\lvert\theta\rvert\,;
\end{equation}
that the slope is constant between the walls is the boundary condition of
the construction, and the fibre's smooth backreaction, absorbed in $V$,
alters the shape between the walls without touching what follows. At the
corner, integrating \eqref{eq:bps} across $[-\epsilon,\epsilon]$,
\begin{equation}\label{eq:kink}
  [\partial_{\theta}a]_{0}\;=\;-\frac{\kappa_{2}^{2}}{2}\int_{-\epsilon}^{\epsilon}\frac{\lvert F_{\psi\omega}\rvert}{\sqrt{g_{2}}}\,\mathrm{d}\theta
  \;=\;-\kappa_{2}^{2}N\pi\,,
\end{equation}
exactly: each unit of flux contributes a kink of exactly $\kappa_{2}^{2}\pi$,
and no smooth term contributes as $\epsilon\to0$. From \eqref{eq:bulk} the
same kink is $a'(0^{+})-a'(0^{-})=-2k$, so
\begin{equation}\label{eq:balance}
  2k\;=\;\kappa_{2}^{2}\,N\pi\,.
\end{equation}
\rigR\ (the jump); \rigD\ (constant slope).
%% Q16: jump exact, slope a boundary condition. Stated as such.

\subsection{Gauss--Bonnet normalisation}\label{sec:gb}

The half-period $L$ and the two-dimensional coupling are not independent.
By Proposition~\ref{prop:chi-pillowcase} the total curvature $4\pi$ sits as
$\pi$ at each corner; the normalisation connecting the deficit angles to
the length of the domain fixes
\begin{equation}\label{eq:kappa-fix}
  \kappa_{2}^{2}\;=\;\frac{2}{L}\,,
\end{equation}
which, with \eqref{eq:kappa-eff}, is a relation among the scales of the
geometry: $\kappa_{12}^{2}=2V_{4}\mathrm{Vol}_{0}e^{6a_{0}}/L$, the
half-period measured in the gravitational unit weighted by the fibre
volume. It is the one normalisation of this section and we tag it as
such. \rigD
%% Assumption 4 of the WP4 list: κ₂² = 2/L is reproduced as in Paper 2 / the
%% Addendum, with eq:kappa-eff making it a scale relation. A derivation from
%% the orbifold Gauss–Bonnet theorem with the e^{6a} weights would upgrade it
%% to [R]; not supplied here. PI to confirm the reading.

\subsection{The result}\label{sec:result}

Substituting \eqref{eq:kappa-fix} in \eqref{eq:balance},
$2k=(2/L)N\pi$, hence $kL=N\pi$, which is Theorem~\ref{thm:warp-exponent}
(\texttt{L5\_balance}). The smallest case by hand: $N=1$, one unit at one
corner, kink $\kappa_{2}^{2}\pi=2\pi/L$, rate $k=\pi/L$, exponent across
the half-period exactly $\pi$; the general case is $N$ copies of this
balance. Corollary~\ref{cor:zeta} is the composition with
Proposition~\ref{prop:flux}: $\zeta=kL=\chi(\Keight)\pi=12\pi$
(\texttt{L5\_zeta}). Three checks: (i) if the fundamental domain
degenerates to $S^{1}/\Ztwo$ the two-form flux becomes a brane tension
and \eqref{eq:kink} is the Israel junction condition of \cite{RS1999};
(ii) the saturated configuration has positive-definite flux energy;
(iii) $\partial_{\theta}a=-k<0$ between the walls, so the warp decreases
monotonically. \rigR

\subsection{Which direction, and the two sheets}\label{sec:direction}

Nothing above depends on whether the transverse coordinate $\theta$ is
$\theta_{\psi}$ or $\theta_{\omega}$: the normalisation \eqref{eq:kappa-fix}
is applied per direction and \eqref{eq:kL-Npi} is direction-blind. The
sheet exchange $\tau$ of Section~\ref{sec:handedness} fixes $P_{1}$ and
$P_{4}$ and exchanges $P_{2}$ and $P_{3}$, the $\psi$- and
$\omega$-half-periods (\texttt{L5\_tau\_corners}), so a profile transverse
to the wall of one sheet is, under $\tau$, a profile transverse to the wall
of the other. The exponent $12\pi$ is therefore $\tau$-symmetric, and
calling the transverse period $L_{\omega}$ or $L_{\psi}$ is the choice of
which sheet one calls $\psi$, the convention of M1 Appendix~B. What is not
conventional is the aspect ratio: the two half-periods differ by
$\tau_{2}=L_{\omega}/L_{\psi}=1.8466$ (M1 Proposition~8.2), so the rate $k$
differs by $\tau_{2}$ between the sheets, and any quantity built from a
period length rather than from the exponent depends on which sheet is
described. In this series ``$L_{\omega}=12\pi$'' means the half-period
transverse to the $\psi$-wall, in the sheet called $\psi$, in the unit of
\eqref{eq:kappa-fix}. \rigD
%% PI ruling Q15. Register entry 16 (Paper 2 θ_ω vs Addendum θ_ψ).

\subsection{What is derived and what is read}\label{sec:summary4}

The jump \eqref{eq:kink}, the balance \eqref{eq:balance}, the integers of
Section~\ref{sec:integers} and the composition \eqref{eq:zeta} are
theorems (\rigR). Read, and named: the placement of the twelve units on
the base cycle (\rigD, Proposition~\ref{prop:flux}); the constant slope
between walls (\rigD, boundary condition); the normalisation
\eqref{eq:kappa-fix} (\rigD); the tie $a=-\zeta\QCD$ with $\QCD=1$ at the
far wall (\rigD, the definition of the dilaton's normalisation); and the
balance of $V$ as the determination of $\Lambda_{12}$ (\rigD). The value
$\zeta=12\pi$ rests on these and on nothing outside this paper and M1.

%%=====================================================================
%% M2 — WP5: what ω is. SURGICAL FRAGMENT. UK English. "We".
%% PI rulings: Q4 (revised; the two zeros derived in M2, GUE maximum [M],
%% spectral proof to M3/M4), Q3 (half-circle route (a) after R3.1), R1
%% (ω = ψ + ψ̄). Depends on WP0 (def:omega, prop:parity), WP2 (thm:handedness,
%% M1 Thm 4.2 drive), WP4 (eq:a-Phi, decaying warp). Placeholder numbering.
%%=====================================================================

\section{What the Coherent Field is}\label{sec:omega-shape}

The Coherent Field entered in Definition~\ref{def:omega} as the even
combination of the sheets. This section says what it measures. The
programme reads $\omega$ as a one-hump quantity: zero when the sheets are
rigidly locked, zero when they are uncorrelated, positive between. We
derive the two zeros from A2 and A4, show where the maximum of the
coupling sits, settle the sign of $\omega$ on the solutions of
Section~\ref{sec:warp}, and name what this paper does not prove.

\subsection{The phase of a sheet field}\label{sec:phase}

Work in the corner reading of Definition~\ref{def:psi}, where each
component is complex, and write $\psi=re^{i\vartheta}$ in $\mathbb{C}_{u}$,
so that $\omega=2r\cos\vartheta$ and $\psi-\bar\psi=2ir\sin\vartheta$; the
statements below hold componentwise, and the module trace sums them. The phase
$\vartheta$ is the angle between the two sheets: $\vartheta=0$ when they
coincide, $\vartheta=\pm\pi/2$ when growth stands exactly against decay.
By Definition~\ref{def:psi}, $\hat\tau$ sends $\vartheta\mapsto-\vartheta$.

\begin{proposition}[the locked zero]\label{prop:locked}
$\omega=0$ at a point iff $\bar\psi=-\psi$ there, i.e. iff the sheet field
is $\hat\tau$-odd, $\vartheta=\pm\pi/2$. At such a point the dilaton
vanishes, $\QCD=0$, and the geometry does not respond to the sheets.
\end{proposition}

\begin{proof}
The Coherent Field is $\psi$ plus its shadow. It is zero at a point exactly
when the shadow there is the negative of the field, $\hat\tau\psi=-\psi$:
when exchanging growth and decay reverses the sign of every component. In
the corner reading, where a component is $re^{i\vartheta}$ and the exchange
sends $\vartheta$ to $-\vartheta$, this says $re^{i\vartheta}=-re^{-i\vartheta}$,
that is $\cos\vartheta=0$, so $\vartheta=\pm\pi/2$: the two sheets stand at
right angles, growth exactly against decay. At such a point the dilaton
$\QCD=\omega\lvert\psi-\bar\psi\rvert^{2}$ is zero because its first factor is,
and so the $\zeta$-term of \eqref{eq:master} vanishes there: the geometry
does not respond. \rigR
\end{proof}

This is A2 read on the field: when growth is exactly opposed by decay the
two sheets are locked, nothing in the geometry changes, and the even
combination is zero. It is the reason the Coherent Field is the
combination $\psi+\bar\psi$ and not the modulus, which does not vanish
there (Section~\ref{sec:omega}). We call this the \emph{locked} or rigid
configuration.
%% PI ruling R1, now a proposition.

\subsection{The uncorrelated zero, from the drive}\label{sec:uncorrelated}

The second zero comes from A4. The drive $M_{0}$ acts on the base torus
(M1 Theorem~4.2), and a sheet field is carried along its orbit,
$\psi\mapsto\psi\circ M_{0}^{n}$. We make precise the average that the
programme's ``orbit average'' refers to.

\begin{definition}[orbit average]\label{def:orbit-average}
For a sheet field $\psi$ and a point $z$ of the base, the orbit average of
$\omega$ over $T$ circuits is
$\langle\omega\rangle_{T}(z)=\tfrac1T\sum_{n<T}\omega(M_{0}^{n}z)$. The
\emph{phase distribution} of $\psi$ is the push-forward of Lebesgue
measure on the base under $z\mapsto\vartheta(z)$.
\end{definition}

\begin{proposition}[the uncorrelated zero]\label{prop:uncorrelated}
The drive is ergodic for Lebesgue measure, so for almost every $z$,
$\langle\omega\rangle_{T}(z)\to\int\omega\,\mathrm{d}\mu$ as $T\to\infty$:
the orbit average is the space average. If the phase distribution of
$\psi$ is uniform on the circle --- the phase of the sheets is
uncorrelated with position --- then $\int\omega\,\mathrm{d}\mu=0$.
\end{proposition}

\begin{proof}
A hyperbolic toral automorphism is ergodic, indeed mixing, for Lebesgue
measure \cite{KatokHasselblatt}; M1 Proposition~4.1 shows A4 forces
hyperbolicity. Birkhoff's theorem gives the first statement. For the
second, with $r$ bounded and $\vartheta$ uniform,
$\int\omega=2\int r\cos\vartheta\,\mathrm{d}\mu=2\bar r\int_{0}^{2\pi}\cos\vartheta\,\tfrac{\mathrm{d}\vartheta}{2\pi}=0$,
where the factorisation uses that a uniform phase is independent of the
modulus. \rigR\ (given the independence, which is what ``uniform phase''
means).
\end{proof}

This is A4 read on the field: a drive of positive entropy mixes, and a
sheet field whose phase has been mixed to uniformity has no net even
combination. We call this the \emph{uncorrelated} configuration.

\subsection{The hump, and where the coupling is largest}\label{sec:hump}

Between the two zeros the average is positive. Let the phase distribution
be concentrated around a lock angle $\vartheta_{0}$ with concentration
$c\ge0$ (any unimodal family will do; the circular normal
$\propto e^{c\cos(\vartheta-\vartheta_{0})}$ is the standard one). Then
\begin{equation}\label{eq:hump}
  \langle\omega\rangle\;=\;2\bar r\,A(c)\cos\vartheta_{0}\,,\qquad
  A(c)=\frac{I_{1}(c)}{I_{0}(c)}\in[0,1)\,,
\end{equation}
which vanishes at $c=0$ (uniform, Proposition~\ref{prop:uncorrelated}),
vanishes at $\vartheta_{0}=\pm\pi/2$ (locked,
Proposition~\ref{prop:locked}), and is positive for
$\lvert\vartheta_{0}\rvert<\pi/2$: one hump in the lock angle between the
two locked configurations. \rigD

The coupling to curvature is not maximal where $\omega$ is. In the phase,
\begin{equation}\label{eq:Phi-phase}
  \QCD\;=\;\omega\,\lvert\psi-\bar\psi\rvert^{2}\;=\;8r^{3}\cos\vartheta\,\sin^{2}\vartheta\,,
\end{equation}
which vanishes at $\vartheta=0$ (sheets coincident, no held distinction)
and at $\vartheta=\pm\pi/2$ (locked), and is extremal where
$\tan^{2}\vartheta=2$, $\sin^{2}\vartheta=\tfrac23$. The geometry responds
most strongly neither when the sheets agree nor when they oppose, but at
the definite intermediate angle $\vartheta^{\ast}=\arctan\sqrt2\approx54.7^{\circ}$
(\texttt{L6\_Phi\_extremum}). \rigR
%% Direct computation from prop:parity; the angle is the parity theorem's
%% own number. Not in the corpus as far as read; flagged as new.

\subsection{The sign of $\omega$ on the solutions}\label{sec:half-circle}

Section~\ref{sec:warp} gave the warp as a decaying exponential,
$a=-\zeta\QCD$ with $\QCD$ rising from $0$ at the wall to $1$ at the far
wall, $\QCD$ there standing for the real part of the trace. On those
solutions, therefore, $\QCD\ge0$. But $\QCD$ is $\omega$ times a
non-negative real number, $\lvert\psi-\bar\psi\rvert^{2}$, so wherever the
sheets are distinct --- wherever that number is not zero --- the real part
of $\omega$ has the same sign as $\QCD$ and is non-negative: in the corner
reading, $\cos\vartheta\ge0$, which is the half-circle
$\lvert\vartheta\rvert\le\pi/2$. The Coherent Field does not change sign on
the solutions. What fixed the sign was the direction of the warp: the
exponential decays away from the wall, so $\QCD$ rises. Had the warp grown
instead, the half-circle would be the other one, and the overall sign of
$\psi$ --- which no axiom fixes, since $\psi\mapsto-\psi$ commutes with
$\hat\tau$ and reverses $\omega$ and $\QCD$ together --- would be chosen to
restore $\mathrm{Re}\,\omega\ge0$. Away from the corner reading $\omega$
carries an imaginary part along the selected unit, which the sign statement
does not constrain. \rigD
%% Closes R2.4 (Q3, route (a)). PI ruling 9 Nov 2026: modulus squared; the
%% signs of Φ and Re ω now agree directly. Minimality Register M2-9 updated.

\subsection{What this paper does not prove}\label{sec:omega-hook}

The programme's canonical reading of $\omega$ is a one-hump complexity
functional of a \emph{spectrum}: zero for a rigidly ordered spectrum,
zero for an uncorrelated one, maximal in the regime of strong level
repulsion, with spectral rigidity as the computable proxy. The two zeros
derived above are the field-level counterparts of the first two
statements: locked for rigidly ordered, uniform for uncorrelated. The
identification of the field $\omega$ with that spectral functional, and
the placement of its maximum in the regime of strong level repulsion, are
not proved here; they require the spectrum of the Dirac operator at the
cone points of the pillowcase and a bridge theorem between the two
descriptions, which is the subject of a later paper of the series. We
record the reading as motivated and the bridge as open. \rigM
%% Q4 (revised): GUE maximum [M]; spectral proof registered as M3/M4
%% (Register §12.1). ω named per omega-definition-usage §6a throughout.


\section{Applications and outlook}\label{sec:outlook}

The equation \eqref{eq:master} with $\zeta=12\pi$ is the one the papers of the programme have used, and this paper fixes its origin: the form from the variation, the coupling from the warp exponent, the variables from the sheets. By axiom: closure gives the equation and the vacuum term, order gives matter and the forces, growth and decay gives the additional term, the drive gives the arrow. Three things follow for the physics without further input. The two sectors $\psi$ and $\bar\psi$ are distinct by handedness relative to the flow, so the shadow sector of the programme is not a relabelling. The stress tensor of the dilaton along the flow is stiff, so wherever $\QCD$ is driven the geometry responds as the programme's collapse and budget analyses assume, read along the flow. And the warp falls by exactly $12\pi$ across a half-period of the base, which is the number from which the programme's hierarchies descend. The papers that follow derive from the same second axiom, through the walls of the sheets and the phase of the dilaton, the masses, the three generations and the mixing angles of the quarks and leptons and the Higgs sector: the structure of matter is the structure of growth and decay, and what a theory of four forces left unexplained, why there is matter and why it does not run backwards, is supplied by the two axioms such a theory omits. We do not pursue those derivations here. \rigM

Four things are left open and named. The spectral face of $\omega$, its identification with a one-hump functional of the Dirac spectrum at the cone points, is the subject of a later paper (Section~\ref{sec:omega-hook}). The normalisation $\kappa_{2}^{2}=2/L$ is reproduced here as a relation among the scales and tagged \rigD; a derivation from the orbifold Gauss--Bonnet theorem with the warped weights would complete it (Section~\ref{sec:gb}). The content of the matter stress tensor, its generations and masses, belongs to the papers that follow (Section~\ref{sec:matter}). And the ratio of the drive's rate per circuit, $\log\varphi^{2}$, to the warp's exponent per half-period, $12\pi$, is a pure number of the geometry whose meaning this paper does not determine (Section~\ref{sec:handedness}).

\section*{Acknowledgements}
The author used Claude (Anthropic) and Gemini (Google DeepMind) as mathematical development and analytical tools during the preparation of this work. All intellectual content, theoretical framework, and scientific claims are the sole responsibility of the author.

\appendix

\section{Machine-verified certificates}\label{app:lean}

Seven certificates accompany the paper, each a single Lean~4.15.0 file on the
core toolchain with no Mathlib dependency, no author axioms and no
\texttt{sorry}. They build within the series' proof tree
\texttt{MEFProof.Paper\_M2}; the build record is
\begin{lstlisting}
Built MEFProof.Paper_M2.M2_L1_parity      (1.2s)
Built MEFProof.Paper_M2.M2_L2_master      (549ms)
Built MEFProof.Paper_M2.M2_L3_zeta        (626ms)
Built MEFProof.Paper_M2.M2_L4_handedness  (1.2s)
Built MEFProof.Paper_M2.M2_L5_warp        (649ms)
Built MEFProof.Paper_M2.M2_L6_omega       (608ms)
Built MEFProof.Paper_M2.M2_L7_module      (863ms)
\end{lstlisting}
Every certified statement is named at the point in the text where it is
used. The analytic steps --- the variation, the Bianchi contraction, the
ergodic theorem, the Pin anticommutation --- are one-line computations or
cited theorems and are not certified; the certificates fix the arithmetic
and the finite case analyses those steps rely on.

\texttt{M2\_L1\_parity.lean} (Section~\ref{sec:variables}, 12 theorems):
the $\hat\tau$-sign of a monomial $(\mathrm{Re}\,\psi)^{a}(\mathrm{Im}\,\psi)^{b}$
is $(-1)^{b}$; vanishing on $\psi=\bar\psi$ iff $b\ge1$ and on $\psi=-\bar\psi$
iff $a\ge1$; degree bound $a+b\ge3$ and uniqueness of $(1,2)$; the parity
table of $\omega$, $\psi-\bar\psi$, $\lvert\psi-\bar\psi\rvert^{2}$, $\QCD$;
multiplicativity of the sign.

\texttt{M2\_L2\_master.lean} (Section~\ref{sec:equation}, 6): the flow
components of $T^{\QCD}$ and the trace in twelve dimensions; $\lvert\nabla\QCD\rvert^{2}$
is $\hat\tau$-even of degree six; the competing two-derivative terms fail
the vanishing conditions.

\texttt{M2\_L3\_zeta.lean} (Section~\ref{sec:warp}, 7): $\chi(\Keight)=3\cdot2\cdot2$;
deficit $\pi$ per cone, total $4\pi=2\pi\chi$; $\zeta/\pi=12$ and the
identity of the cone-deficit and Euler readings.

\texttt{M2\_L4\_handedness.lean} (Section~\ref{sec:handedness}, 11):
$SM_{0}S^{-1}=M_{0}^{-1}$; $S^{2}=-I$; $\det S=1$ and the orientation of
$\tau$; the eigen-directions $(\varphi,1)$, $(-1,\varphi)$ in $\Z[\varphi]$
with eigenvalues $\varphi^{\pm2}$; $S$ exchanges them; any involutive flow
reversal has determinant $-1$.

\texttt{M2\_L5\_warp.lean} (Section~\ref{sec:warp}, 9): the orbifold Euler
count; the fundamental-domain check; the balance $kL=N\pi$; $\zeta=12\pi$
and the $N=1$ case; $\tau$ fixes $P_{1},P_{4}$ and exchanges $P_{2},P_{3}$.

\texttt{M2\_L6\_omega.lean} (Section~\ref{sec:omega-shape}, 5): the
coefficient $8$ in $\QCD=8r^{3}\cos\vartheta\sin^{2}\vartheta$; the
extremum $\tan^{2}\vartheta=2$; the sign bookkeeping of the half-circle; the
zeros of the hump.

\texttt{M2\_L7\_module.lean} (Sections~\ref{sec:variables}, \ref{sec:equation}, 8):
$\mathrm{Cl}(0,12)\cong\Hquat(32)$ by Bott periodicity ($12\equiv4\bmod8$) and
the module dimensions $32$, $64$, $128$; the three readings of the kinetic
term, $1$, $32$, $1024$; the trace of a module-symmetric vacuum; the square
of an imaginary quaternion is minus its squared modulus.

The seven files follow verbatim.

\subsection*{\texttt{M2\_L1\_parity.lean} --- the parity algebra}
\lstinputlisting{../../MEFProof/Paper_M2/M2_L1_parity.lean}

\subsection*{\texttt{M2\_L2\_master.lean} --- the dilaton term}
\lstinputlisting{../../MEFProof/Paper_M2/M2_L2_master.lean}

\subsection*{\texttt{M2\_L3\_zeta.lean} --- the integers of $\zeta$}
\lstinputlisting{../../MEFProof/Paper_M2/M2_L3_zeta.lean}

\subsection*{\texttt{M2\_L4\_handedness.lean} --- handedness}
\lstinputlisting{../../MEFProof/Paper_M2/M2_L4_handedness.lean}

\subsection*{\texttt{M2\_L5\_warp.lean} --- the warp exponent}
\lstinputlisting{../../MEFProof/Paper_M2/M2_L5_warp.lean}

\subsection*{\texttt{M2\_L6\_omega.lean} --- the shape of $\omega$}
\lstinputlisting{../../MEFProof/Paper_M2/M2_L6_omega.lean}

\subsection*{\texttt{M2\_L7\_module.lean} --- the Clifford module}
\lstinputlisting{../../MEFProof/Paper_M2/M2_L7_module.lean}

%%=====================================================================
%% Appendix B seed — cumulative symbol table (M1 reproduced, M2 extended).
%% To be placed as M2 Appendix B; the M1 block is verbatim from M1 v2 App. B.
%%=====================================================================

\section{The labelling of the corners, and the symbols of the series}\label{app:labels}

\paragraph{From M1 (reproduced).}
The four corners are named by position, $P_{1}=(0,0)$, $P_{2}=(\tfrac12,0)$,
$P_{3}=(0,\tfrac12)$, $P_{4}=(\tfrac12,\tfrac12)$ in coordinates
$(\theta_{\psi},\theta_{\omega})$, and by algebra, $C_{0}=\R$, $C_{i}$,
$C_{j}$, $C_{k}$. Throughout the programme the dictionary is: $P_{1}=C_{0}$
(fixed by the drive; $\iota$ acts as conjugation); $P_{3}=C_{i}$ (the
$\omega$ half-period); $P_{2}=C_{j}$ (the $\psi$ half-period); $P_{4}=C_{k}$
(the corner at which the $\Spinc$ lift carries the sign $-1$). The assignment
of $j$ and $k$ between $P_{2}$ and $P_{3}$ is a convention and carries no
content (M1 Appendix~B).

\paragraph{Introduced in M2.}
\begin{center}\small
\begin{tabular}{@{}p{0.29\textwidth}p{0.49\textwidth}p{0.15\textwidth}@{}}
\toprule
Symbol & Meaning & Where defined \\
\midrule
$\psi,\ \bar\psi=\hat\tau\psi$ & the sheet field, a section of the Clifford module $M\cong\Hquat^{32}$, and its shadow & Def.~\ref{def:psi} \\
$M$, $\langle\cdot,\cdot\rangle_{M}$ & irreducible $\mathrm{Cl}(0,12)$-module; its inner product (trace) & Def.~\ref{def:psi}, \S\ref{sec:action} \\
$\mathcal{A},\ \mathcal{C}$ & bundle of fibre algebras; corner-algebra subbundle (the corner reading) & Def.~\ref{def:psi} \\
$\tau,\ \hat\tau$ & flow-reversing involution of A2 ($\tau^{2}=$ fold on the cover); its lift & \S\ref{sec:handedness} \\
$S$, $v_{u}$, $v_{s}$, $\kappa=\log\varphi^{2}$ & quarter-turn; unstable/stable directions of $M_{0}$; Sol rate & \S\ref{sec:handedness} \\
$\omega=\psi+\bar\psi$ & the Coherent Field; $\hat\tau$-even, quaternion-valued, carries the imaginary direction & Def.~\ref{def:omega} \\
$\psi-\bar\psi$ & the odd sheet variable; $\hat\tau$-odd & Def.~\ref{def:omega} \\
$\QCD=\omega\lvert\psi-\bar\psi\rvert^{2}$ & the Quantum Coherence Dilaton, componentwise in $M$ & Prop.~\ref{prop:parity} \\
$A$, $w=e^{A}$ & the A2 warp exponent and warp factor on the base & Def.~\ref{def:scheme} \\
$k_{\omega}$ & decay rate of $A$ along $\theta_{\omega}$ & Section~\ref{sec:warp} \\
$F$ & curvature of the $\Spinc$ determinant line & Def.~\ref{def:scheme} \\
$\kappa_{2}$ & two-dimensional gravitational coupling on the base & Section~\ref{sec:warp} \\
$T^{\QCD}_{MN}$ & stress tensor of $\QCD$, the module trace of the scalar form & Section~\ref{sec:equation} \\
$\zeta$ & coupling of $\QCD$: coefficient of $T^{\QCD}_{MN}$; equals the warp exponent per half-period & Sections~\ref{sec:equation}--\ref{sec:warp} \\
$G_{\mathrm{eff}},\ \Lambda_{\mathrm{eff}}$ & effective gravitational and vacuum couplings & Section~\ref{sec:equation} \\
\bottomrule
\end{tabular}
\end{center}
Later papers of the series reproduce this table and extend it.


\begin{thebibliography}{99}
\bibitem{BishopONeill1969} R.~L. Bishop and B. O'Neill, Manifolds of negative curvature, \emph{Trans. Amer. Math. Soc.} 145 (1969) 1--49.
\bibitem{Cartan1922} \'E. Cartan, Sur les \'equations de la gravitation d'Einstein, \emph{J. Math. Pures Appl.} 1 (1922) 141--203.
\bibitem{Friedrich} T. Friedrich, \emph{Dirac Operators in Riemannian Geometry}, AMS, 2000.
\bibitem{KatokHasselblatt} A. Katok and B. Hasselblatt, \emph{Introduction to the Modern Theory of Dynamical Systems}, Cambridge University Press, 1995.
\bibitem{Kawasaki1979} T. Kawasaki, The Riemann--Roch theorem for complex $V$-manifolds, \emph{Osaka J. Math.} 16 (1979) 151--159.
\bibitem{LawsonMichelsohn} H.~B. Lawson and M.-L. Michelsohn, \emph{Spin Geometry}, Princeton University Press, 1989.
\bibitem{Lovelock1971} D. Lovelock, The Einstein tensor and its generalizations, \emph{J. Math. Phys.} 12 (1971) 498--501.
\bibitem{ONeill1983} B. O'Neill, \emph{Semi-Riemannian Geometry}, Academic Press, 1983.
\bibitem{RS1999} L. Randall and R. Sundrum, A large mass hierarchy from a small extra dimension, \emph{Phys. Rev. Lett.} 83 (1999) 3370--3373.
\bibitem{Thurston} W.~P. Thurston, \emph{The Geometry and Topology of Three-Manifolds}, Princeton lecture notes, 1980.
\bibitem{Vermeil1917} H. Vermeil, Notiz \"uber das mittlere Kr\"ummungsma\ss\ einer $n$-fach ausgedehnten Riemann'schen Mannigfaltigkeit, \emph{Nachr. Ges. Wiss. G\"ottingen} (1917) 334--344.
\bibitem{Wald} R.~M. Wald, \emph{General Relativity}, University of Chicago Press, 1984.
\bibitem{Weyl1922} H. Weyl, \emph{Raum, Zeit, Materie}, 5th ed., Springer, 1923.
\end{thebibliography}
\end{document}
